\documentclass{ieeeaccess}

\usepackage{cite}
\usepackage{amsmath,amssymb}
\usepackage{booktabs}
\usepackage{multirow}
\usepackage{graphicx}
\usepackage{xspace}
\usepackage{float}
\usepackage{pifont}
\usepackage{array}
\usepackage{xcolor}
\usepackage{url}
\usepackage{hyperref}
\usepackage{algorithm}
\usepackage{algorithmic}
\usepackage{subcaption}
\usepackage{siunitx}
\usepackage{bm}
\usepackage{needspace}
\usepackage{placeins}

\setcounter{topnumber}{4}
\setcounter{bottomnumber}{2}
\setcounter{totalnumber}{6}
\setcounter{dbltopnumber}{3}
\renewcommand{\topfraction}{0.95}
\renewcommand{\bottomfraction}{0.90}
\renewcommand{\textfraction}{0.05}
\renewcommand{\floatpagefraction}{0.80}
\renewcommand{\dbltopfraction}{0.95}
\renewcommand{\dblfloatpagefraction}{0.80}

\newcommand{\model}{PatchFusionBERT\xspace}
\newcommand{\modelvzero}{PatchFusionBERT$_{\text{v0}}$\xspace}
\newcommand{\modelvtwo}{PatchFusionBERT$_{\text{v2}}$\xspace}
\newcommand{\patchtst}{PatchTST\xspace}
\newcommand{\dlinear}{DLinear\xspace}
\newcommand{\bertonly}{BERTOnly\xspace}
\newcommand{\cmark}{\ding{51}}
\newcommand{\xmark}{\ding{55}}
\newcommand{\safeincludegraphics}[2][]{%
  \IfFileExists{#2}{\includegraphics[#1]{#2}}{}
}

\hypersetup{
    colorlinks=true,
    linkcolor=blue,
    citecolor=blue,
    urlcolor=blue,
    pdftitle={PatchFusionBERT: Post-Encoder Fusion for Patch-Based Long-Horizon Time Series Forecasting},
    pdfauthor={Hussein Ahmad}
}

\begin{document}

\vol{4}
\year{2026}
\history{}
\doi{}

\title{PatchFusionBERT: Post-Encoder Fusion for Patch-Based Long-Horizon Time Series Forecasting}

\author{Hussein Ahmad\authorrefmark{1}*,
Seyyed Kasra Mortazavi\authorrefmark{1},
Taha Benarbia\authorrefmark{2},
Fadi Al Machot\authorrefmark{3}, and
Kyandoghere Kyamakya\authorrefmark{1}}
\address[1]{Institute for Smart Systems Technologies, Universität Klagenfurt, Klagenfurt, Austria}
\address[2]{Industrial Engineering Department, University of Oran 2, Oran, Algeria}
\address[3]{Department of Data Science, Norwegian University of Life Sciences, Ås, 1432, Norway}

\markboth
{Ahmad \headeretal: PatchFusionBERT}
{Ahmad \headeretal: PatchFusionBERT}

\corresp{Corresponding authors: Hussein Ahmad (e-mail: hussahmad@edu.aau.at) and Kyandoghere Kyamakya (e-mail: kyandoghere.kyamakya@aau.at).}

% =========================
\begin{abstract}
Long-horizon time-series forecasting requires architectures that
capture extended temporal structure while remaining computationally
tractable.
We propose \model, a patch-based forecasting architecture that
augments a PatchTST-style backbone with a compact post-encoder
Transformer refinement stage and explicit fusion between backbone
and refined patch representations.
Two fusion variants are studied under a unified pipeline across
seven public benchmarks (21 dataset--horizon settings) and a
parameter-matched PatchTST control that isolates the fusion
contribution from raw model capacity: direct concatenation
(\modelvzero) and concatenation followed by learnable projection
(\modelvtwo).
This study covers supervised point forecasting. It does not address
probabilistic forecasting, transfer learning, or foundation-model
pretraining.
Under the 3-run median capacity-control protocol, both fusion variants achieve
lower MSE than width-scaled PatchTST on Weather and ETTh2 at every tested
horizon, providing the strongest evidence for an effect not explained by width
scaling alone.
ETTm2 is partially favorable, whereas Exchange and ETTm1 are
neutral-to-negative, indicating that the gains are
regime-dependent rather than universal.
On Illness, the apparent fusion advantage observed in the
single-seed benchmark does not survive the aligned multi-seed
capacity-matched follow-up, where PatchTST variants achieve
lower mean MSE at all three horizons. We therefore reinterpret
Illness as evidence for added patch-based capacity in a low-data
regime rather than for fusion specifically.
A single-seed K-depth ablation on ETTh2 and Weather is consistent
with most of the benefit being captured at $K{=}1$ in the tested
setting, with performance largely insensitive to $K$ within
$\{1,2,3\}$.
Overall, post-encoder refinement combined with explicit fusion is
a data-dependent design lever for patch-based long-horizon
forecasting, not a universal architectural improvement.
Additional analyses --- five-seed stability evaluation, pooled Wilcoxon signed-rank
tests, head-capacity and depth-matched controls, broader robustness probes, CKA
diagnostics, and dataset-regime characterization --- consistently support this
bounded interpretation without changing the main conclusions.
\end{abstract}

\begin{keywords}
time series forecasting, long-term forecasting, transformers, patch-based modeling,
post-encoder fusion, representation learning, multi-horizon prediction
\end{keywords}

\titlepgskip=-15pt

\maketitle

% =========================
\section{Introduction}
% =========================

Time-series forecasting is a foundational machine-learning problem because many real
systems make decisions from evolving observations rather than static inputs. Energy
demand management, traffic control, weather and climate monitoring, public-health
surveillance, and financial analytics all rely on forecasts that remain useful not
only one step ahead, but across longer future windows where uncertainty compounds and
temporal dependencies are harder to exploit. In long-horizon settings, practical
forecasting models must balance three competing requirements: preserving sufficient
historical context, remaining computationally tractable as sequence length grows, and
generalizing across datasets with different sampling rates, noise levels, and
cross-variate structure \cite{zhou2021informer,wu2021autoformer,kim2022revin}.

Transformer-based methods became a major line of work for this setting because
self-attention offers a flexible mechanism for modeling long-range dependencies
\cite{vaswani2017attention}. Early long-sequence forecasters such as Informer and
Autoformer pursued this direction with redesigned attention and decomposition
mechanisms \cite{zhou2021informer,wu2021autoformer}. At the same time, the field has
become more critical and more precise: stronger evidence shows that architectural
complexity does not automatically translate into better forecasting, and linear or
minimalist baselines can remain highly competitive on standard benchmarks
\cite{zeng2023transformers}. This has shifted emphasis toward identifying which
components help, in which regimes, and at what computational cost.

One of the clearest advances in this design space is the move from point-wise
tokenization to patch-based representations. PatchTST showed that grouping consecutive
timesteps into patches reduces effective sequence length while preserving local temporal
semantics \cite{nie2023patchtst}. Subsequent models explored complementary
representation choices: iTransformer emphasizes variate tokens and cross-variate
dependency learning \cite{liu2024itransformer}, TimeMixer advances multi-scale mixing
strategies \cite{wang2024timemixer}, TiDE proposes a dense encoder-decoder with
multi-layer perceptrons \cite{das2023tide}, and TimeXer highlights
endogenous-exogenous integration in practical forecasting settings
\cite{wang2024timexer}. Recent analyses further question whether self-attention itself
is always the dominant source of gains, stressing the roles of representation choice
and architectural inductive bias \cite{selfattention2024tsf,closerlook2024tsf}.

In parallel, foundation-model directions such as Chronos and broader pattern-learning
frameworks such as TimeMixer++ suggest that larger-scale pretraining and more universal
time-series representations can improve cross-domain transfer
\cite{ansari2024chronos,timemixerpp2025}. These developments are important, but they
do not remove the need for controlled supervised benchmark studies. In many practical
research settings, the relevant question is whether a compact architecture can improve
accuracy in standard forecasting pipelines while staying interpretable, reproducible,
and computationally reasonable.

Within this context, patch-based forecasting offers a strong efficiency-accuracy
starting point, but it leaves an open architectural question. Patching shortens token
sequences and structures local history, yet a shallow patch encoder may under-refine
useful interactions at longer horizons or under heterogeneous temporal dynamics. In
NLP, deep encoder stacks such as BERT popularized repeated contextual refinement
\cite{devlin2019bert}. That idea should not be transferred naively to forecasting,
but it motivates a forecasting-specific question: can a compact post-encoder
refinement stage improve patch representations without discarding the computational
advantages of patching?

\textbf{Contributions.} Against this background, this paper makes the following
contributions.
\begin{enumerate}
\item We introduce \model, which combines patch tokenization, Transformer-style
post-encoder refinement, and explicit fusion between backbone and refined
representations, and study two fusion variants: direct concatenation (\modelvzero) and
concatenation with learnable projection (\modelvtwo).
\item We evaluate the approach across seven established long-horizon benchmarks
(21 locked dataset--horizon settings), identifying where post-encoder fusion is
associated with lower error and where simpler or alternative models remain competitive.
\item We conduct multi-axis capacity-matched controls: a width-matched PatchTST
baseline across all horizons, a depth-matched PatchTST variant at $H=192$, and a
head-capacity LargeHead control, to isolate architectural effects from parameter growth.
\item We perform five-seed stability evaluation and pooled Wilcoxon signed-rank tests
(Holm--Bonferroni and BH-FDR correction) to provide global supporting evidence
beyond single-run observations.
\item We include a K-depth ablation, controlled efficiency benchmark, broader
robustness probes (random and block missingness, Gaussian noise, channel dropout),
CKA representation diagnostics, and dataset-regime characterization to build
a coherent mechanistic picture.
\end{enumerate}
A preliminary version of this work appeared as a conference paper
\cite{ahmad2024icamcs_patchbert}. Full details of the technical differences from that
version are given in
Section~\ref{sec:related}.

Taken together, these contributions frame \model\ as an accuracy-oriented extension
of patch-based forecasting rather than a claim of universal architectural superiority.
The following sections proceed from architecture to protocol to evidence.
Section~\ref{sec:methodology} defines the model and its fusion variants.
Section~\ref{sec:experimental_setup} describes the benchmark protocol.
Section~\ref{sec:results} evaluates performance and limitations.

% =========================
\section{Related Work}\label{sec:related}
% =========================

\subsection{Long-Horizon Forecasting and the Reassessment of Model Complexity}
Long-horizon time-series forecasting is a demanding testbed for both representation
quality and scalability. Early Transformer-based forecasters such as Informer and
Autoformer argued that attention-centric designs could improve long-range temporal
modeling through sparse attention and decomposition-based structure
\cite{zhou2021informer,wu2021autoformer}. These works established Transformers as
serious forecasting models, but they also motivated a broader methodological question:
whether architectural sophistication is itself the primary source of improvement, or
whether gains also depend on protocol choices such as tokenization, normalization, and
training setup. This concern became explicit as later work showed that comparatively
simple baselines can remain highly competitive on standard LTSF benchmarks
\cite{zeng2023transformers,closerlook2024tsf}. Our work is motivated by a narrower
and more testable question. The goal is not to claim that deeper models are always
better. Instead, we test whether a specific post-encoder refinement stage adds value
on top of a strong patch-based backbone.

\subsection{Patch-Based and Alternative Tokenization Strategies}
Among efficient forecasters, patch-based modeling has been especially influential.
PatchTST showed that segmenting a series into subseries-level patches shortens
effective token sequences while preserving local temporal semantics
\cite{nie2023patchtst}. Subsequent models explored different interaction structures:
iTransformer inverts the usual view by treating variates as tokens to emphasize
cross-variate dependency learning \cite{liu2024itransformer}, TimeMixer uses
decomposable multiscale mixing to achieve strong performance without standard
self-attention \cite{wang2024timemixer}, TiDE uses a dense encoder-decoder to model
temporal dynamics without explicit attention \cite{das2023tide}, and TimeXer
emphasizes forecasting with exogenous variables, a practically important setting
often underrepresented in benchmark-focused papers \cite{wang2024timexer}. Together,
these studies suggest that forecasting quality depends not only on depth, but on how
temporal structure, scale information, and variable interactions are represented and
combined.

\subsection{Deeper Encoder Refinement and BERT-Inspired Design}
In NLP, BERT popularized deep encoder stacks for iterative contextual refinement
\cite{devlin2019bert}. The key principle is repeated reweighting of token interactions
across layers rather than a single shallow contextualization pass. This principle
cannot be transferred mechanically to forecasting because the objective, causality
constraints, and evaluation setting are different. Still, it motivates a relevant
architectural hypothesis: representations useful after a shallow encoder may still
benefit from additional structured refinement. Related time-series work on
representation quality, including TimeMAE, supports the broader importance of
refinement and representation learning choices \cite{cheng2023timemae,selfattention2024tsf}.
Our approach adopts only the encoder-side refinement idea in a supervised forecasting
pipeline and does not rely on large-scale self-supervised pretraining.

\subsection{Recent Movement Toward Broader Time-Series Foundation Models}
The literature has recently expanded toward broader time-series foundation modeling.
Chronos demonstrates that large-scale pretraining can improve cross-domain transfer
and zero-shot behavior \cite{ansari2024chronos}. More recent frameworks such as
TimeMixer++ also pursue more general pattern-learning formulations across time-series
tasks \cite{timemixerpp2025}. These developments are important, but they do not
eliminate the value of controlled supervised architectural studies. For
benchmark-driven questions under fixed train/validation/test splits, compact models
remain useful because their behavior and trade-offs are often easier to interpret than
those of very large pretrained systems.

\subsection{Position of the Present Work}
Against this landscape, \model\ occupies a specific middle position. It is not a new
long-sequence attention mechanism in the style of Informer, not a decomposition
architecture in the style of Autoformer, not a variate-token reformulation like
iTransformer, not a multiscale mixer in the style of TimeMixer, and not a
dense-encoder design like TiDE
\cite{zhou2021informer,wu2021autoformer,liu2024itransformer,wang2024timemixer,das2023tide}.
Its closest conceptual ancestor is PatchTST, from which it inherits patch tokenization
and efficient backbone processing \cite{nie2023patchtst}. The central difference is an
explicit post-encoder refinement-and-fusion stage that combines backbone features with
refined features before prediction. The purpose is not to assert universal superiority,
but to test whether moderate additional refinement can improve forecasting in selected
regimes while retaining the practical strengths of patch-based processing.

\subsubsection*{Differences from the conference version}
A preliminary version of this work appeared as \cite{ahmad2024icamcs_patchbert},
which introduced the initial PatchTST+BERT integration concept on a limited set of
benchmarks. The present paper extends that work in five substantive dimensions:
(i)~two distinct fusion variants (v0: Concat; v2: Concat+Proj) are formalized and
their trade-offs studied; (ii)~TiDE and TimeXer are added as competitive baselines;
(iii)~a capacity-matched PatchTST control is introduced to separate architectural
from capacity effects; (iv)~a K-depth ablation and multi-seed stability analyses
provide evidence beyond the single-seed benchmark; (v)~the evaluation is expanded
to 21 dataset--horizon settings with robustness and sensitivity studies.

% =========================
\section{Methodology}\label{sec:methodology}
% =========================

This section formalizes \model\ as a supervised patch-based forecaster designed to
answer a focused architectural question: can a compact post-encoder refinement stage
improve patch representations beyond a shallow backbone alone? To isolate this
question, we keep the core PatchTST-style tokenization and backbone processing fixed,
and vary only the refinement-and-fusion stage before the prediction head.

\subsection{Problem Formulation}
We consider the standard supervised multivariate long-horizon forecasting setting.
Given an input sequence $\mathbf{X} \in \mathbb{R}^{L \times C}$, where $L$ is the
lookback length and $C$ is the number of variables, the objective is to predict the
next $H$ future steps $\hat{\mathbf{Y}} \in \mathbb{R}^{H \times C}$:
\begin{equation}
    \hat{\mathbf{Y}} = f(\mathbf{X}; \theta), \quad \hat{\mathbf{Y}} \in \mathbb{R}^{H \times C},
\end{equation}
where $\theta$ denotes learnable parameters. We study point forecasting under the
standard benchmark protocol, optimize the training objective with MSE, and report both
MSE and MAE during evaluation.

\subsection{Patch Tokenization}
Following the patching paradigm of PatchTST \cite{nie2023patchtst}, each variable
channel is segmented into patches of length $P$ with stride $S$:
\begin{equation}
    \mathbf{x}^{(c)} \rightarrow [\mathbf{x}^{(c)}_1, \mathbf{x}^{(c)}_2, \ldots,
    \mathbf{x}^{(c)}_N], \quad \mathbf{x}^{(c)}_i \in \mathbb{R}^{P},
\end{equation}
where $N = \left\lfloor \frac{L - P}{S} \right\rfloor + 1$.
Each patch is mapped into a $d$-dimensional token via a learnable projection and
positional encoding:
\begin{equation}
    \mathbf{z}^{(c)}_i = \mathbf{W}_{p}\mathbf{x}^{(c)}_i + \mathbf{e}_{\text{pos}}(i),
    \quad \mathbf{W}_{p} \in \mathbb{R}^{d \times P}.
\end{equation}
Patch tokenization reduces the effective sequence length from timestep-level
processing to patch-level processing while retaining local temporal continuity inside
each token. In this work, patching is a controlled baseline inherited from PatchTST
rather than a proposed novelty. The architectural contribution begins after this stage.

\subsection{Backbone Patch Encoder (First-Pass Contextualization)}
Let $\mathbf{Z}^{(c)} = [\mathbf{z}^{(c)}_1, \ldots, \mathbf{z}^{(c)}_N] \in
\mathbb{R}^{N \times d}$. A lightweight Transformer encoder (as in PatchTST) provides
initial patch-level contextualization:
\begin{equation}
    \mathbf{H}^{(0,c)} = \text{Enc}_{\text{backbone}}\left(\mathbf{Z}^{(c)}\right).
\end{equation}
In line with PatchTST's channel-independent design \cite{nie2023patchtst}, the same
encoder is applied per channel, keeping complexity controlled and making
$\mathbf{H}^{(0,c)}$ a clean first-pass representation for subsequent refinement
analysis.

\subsection{Post-Encoder Refinement}
The core module in \model\ applies a deeper stack of Transformer blocks to backbone
patch tokens, following the BERT-inspired depth-as-refinement principle on the encoder
side \cite{devlin2019bert}:
\begin{equation}
    \mathbf{H}^{(k,c)} = \text{Block}\left(\mathbf{H}^{(k-1,c)}\right), \quad k=1,\ldots,K.
\end{equation}
Each block contains multi-head self-attention, a position-wise feed-forward network,
residual connections, and layer normalization. In all main benchmark experiments, we
fix $K=3$ for consistency with the original design configuration. A K-depth ablation
sweep (Section~\ref{sec:kdepth}) then shows that performance is largely insensitive
to $K$ within $\{1,2,3\}$, suggesting $K=1$ as a reasonable default in the tested
setting because it captures most of the observed benefit at lower computational cost. The
goal is to iteratively strengthen cross-patch interactions and hierarchical structure
that may remain underdeveloped after a shallow backbone.

\subsection{Fusion Variants}
Let $\mathbf{H}^{(0,c)} \in \mathbb{R}^{N \times d}$ be the backbone output for
channel $c$, and let $\mathbf{H}^{(K,c)} \in \mathbb{R}^{N \times d}$ be the output
after $K$ refinement blocks. We define variant-specific fusion as follows:
\begin{itemize}
    \item \textbf{\modelvzero\ (Concat):}
    \begin{equation}
        \mathbf{U}^{(c)} = [\mathbf{H}^{(0,c)};\mathbf{H}^{(K,c)}],
    \end{equation}
    where $[\cdot;\cdot]$ denotes feature-dimension concatenation. Therefore,
    $\mathbf{U}^{(c)} \in \mathbb{R}^{N \times 2d}$ and this variant directly exposes
    both first-pass and refined representations to the forecasting head.
    \item \textbf{\modelvtwo\ (Concat+Proj):}
    \begin{equation}
        \mathbf{U}^{(c)} = [\mathbf{H}^{(0,c)};\mathbf{H}^{(K,c)}]\mathbf{W}_{f}^{\top}
        + \mathbf{b}_{f},
    \end{equation}
    with $\mathbf{W}_{f} \in \mathbb{R}^{d \times 2d}$,
    $\mathbf{b}_{f} \in \mathbb{R}^{d}$ (broadcast along token positions), and
    $\mathbf{U}^{(c)} \in \mathbb{R}^{N \times d}$. This learnable projection maps
    fused features back to dimension $d$ before the head.
\end{itemize}
In both variants, fusion is applied strictly after patch tokenization and backbone
encoding, and no future timesteps are exposed during encoding or fusion. The term
\textit{BERT-inspired} refers only to encoder-block depth/refinement, not to
masked-language pretraining. Figure~\ref{fig:architecture} summarizes this
end-to-end data flow.

\begin{figure}[t]
\centering
\includegraphics[width=0.98\linewidth]{figures-300426/Fig1_architecture.png}
\caption{Architecture overview of \model. Input patches are encoded by a
PatchTST-style backbone to produce $\mathbf{H}^{(0,c)}$, refined by a $K$-block
post-encoder Transformer to $\mathbf{H}^{(K,c)}$, then fused via Concat
(\modelvzero) or Concat+Proj (\modelvtwo) before the channel-wise forecasting head.}
\label{fig:architecture}
\end{figure}

\subsection{Forecasting Head}
The fused representation is mapped to the prediction horizon through a channel-wise
linear head. Let $d_u=2d$ for \modelvzero\ and $d_u=d$ for \modelvtwo. For each
channel $c$, we flatten the token matrix and apply:
\begin{equation}
    \hat{\mathbf{y}}^{(c)} = \mathbf{W}_{h}^{(c)}\,\mathrm{vec}\!\left(\mathbf{U}^{(c)}\right)
    + \mathbf{b}_{h}^{(c)},
\end{equation}
where $\mathbf{W}_{h}^{(c)} \in \mathbb{R}^{H \times (N d_u)}$ and
$\mathbf{b}_{h}^{(c)} \in \mathbb{R}^{H}$. The final forecast
$\hat{\mathbf{Y}} \in \mathbb{R}^{H \times C}$ is obtained by concatenating outputs
for all $C$ channels.

\subsection{Architectural Comparison}
Table~\ref{tab:architecture} summarizes the conceptual differences between baselines
and \model. Parameter values are reported as representative reference counts or
ranges for the locked benchmark; exact follow-up capacity-control counts are given
separately in Table~\ref{tab:capmatch}.

\begin{table}[t]
\centering
\caption{Architectural comparison of forecasting models with representative
parameter reference values. $^\dagger$Parameter counts vary with forecast horizon
because the linear prediction head scales with $H$; values marked $^\dagger$ are
exact for the locked-snapshot configuration ($d_{\text{model}}=128$,
$d_{\text{ff}}=512$, $P=16$, $S=8$). Exact follow-up capacity-control counts are
reported separately in Table~\ref{tab:capmatch}.}
\label{tab:architecture}
\small
\resizebox{\columnwidth}{!}{%
\begin{tabular}{@{}lcccc@{}}
\toprule
Model & Tokenization & Backbone Depth & Fusion Module & Parameters \\
\midrule
\dlinear       & None (linear)    & 0        & None                       & 0.07M (H=96)–0.23M (H=336) \\
\patchtst      & Patch ($P$=16)   & 3 layers & None                       & 1.1M (H=96), 2.4M (H=336)  \\
PatchTST$_\text{cap}$ & Patch ($P$=16) & deeper/wider & None           & protocol-dependent; see Table~\ref{tab:capmatch} \\
Pointwise+Refine & None (pointwise) & 0        & 3-layer encoder-only stack & 2.8M  \\
\modelvzero    & Patch ($P$=16)   & 3 layers & 3 layers + Concat          & 2.2M (H=96), 4.8M (H=336)$^\dagger$ \\
\modelvtwo     & Patch ($P$=16)   & 3 layers & 3 layers + Concat+Proj     & 1.9M (H=96), 3.2M (H=336)$^\dagger$ \\
\bottomrule
\end{tabular}
}
\end{table}

% =========================
\section{Experimental Setup}\label{sec:experimental_setup}
% =========================

The experimental design evaluates \model\ under a unified long-horizon forecasting
protocol with a focused objective: determine whether post-encoder refinement and
fusion improve prediction quality beyond a strong patch-based backbone, and identify
regimes where gains are substantial, modest, or absent.

\subsection{Datasets}
We evaluate on seven public datasets standard in the long-term forecasting literature:
\begin{itemize}
    \item \textbf{ETT family (ETTh1, ETTh2, ETTm1, ETTm2):} multivariate
    electricity-transformer temperature series with hourly (ETTh) and 15-minute
    (ETTm) resolutions \cite{zhou2021informer,wu2021autoformer}.
    \item \textbf{Exchange:} daily exchange-rate series across multiple countries,
    often structurally harder for highly specialized deep models \cite{lai2018lstnet}.
    \item \textbf{Weather:} multivariate meteorological series with richer
    cross-variable temporal structure \cite{zhou2021informer,wu2021autoformer}.
    \item \textbf{Illness:} influenza-like illness (ILI) signals representing a
    comparatively lower-data forecasting regime \cite{zhou2021informer}.
\end{itemize}
All datasets follow the standard train/validation/test partitions used in benchmark
implementations \cite{zhou2021informer,wu2021autoformer}.

\subsection{Forecasting Setting and Horizons}
We use the standard supervised long-horizon protocol: given a historical window of
length $L$, the model predicts the next $H$ timesteps.
\begin{itemize}
    \item \textbf{Input length:} $L=336$ for ETT/Exchange/Weather and $L=104$ for
    Illness.
    \item \textbf{Forecast horizons:} $H \in \{96,192,336\}$ for
    ETT/Exchange/Weather and $H \in \{24,48,60\}$ for Illness. We restrict the
    locked benchmark to these 21 settings to keep the comparison aligned with
    stable, fully supported dataset--horizon combinations and to avoid mixing
    longer-horizon configurations that are not consistently available across all
    datasets. Horizon $H=720$ is therefore left outside the present scope. For
    descriptive continuity only, Appendix~\ref{app:full_results} includes a
    supplementary single-seed unified-LR bridging view that extends to $H=720$
    (Figure~\ref{fig:app_bridging_lr}); it is not treated as part of the locked
    benchmark.
    \item \textbf{Coverage:} 21 dataset and horizon settings in total.
\end{itemize}

\begin{quote}
\textbf{Protocol boundary.} This paper studies deterministic supervised
point forecasting under fixed public benchmark splits. It does not address
probabilistic forecasting, interval calibration, foundation-model pretraining,
transfer learning, or per-baseline hyperparameter-tuned leaderboard optimization.
\end{quote}

\subsection{Evaluation Metrics}
We report mean squared error (MSE) and mean absolute error (MAE), averaged across
variables (lower is better). The study focuses on point forecasting. Discussion
emphasizes consistent patterns and regime-level differences rather than treating very
small numerical gaps as practically decisive.

\subsection{Baselines}
The baseline set covers representative forecasting families:
\begin{itemize}
    \item \textbf{\dlinear:} a strong low-complexity linear baseline
    \cite{zeng2023transformers}.
    \item \textbf{\patchtst:} the primary architectural reference
    \cite{nie2023patchtst}.
    \item \textbf{iTransformer:} a competitive Transformer baseline with a
    variate-oriented token perspective \cite{liu2024itransformer}.
    \item \textbf{TiDE:} a strong dense encoder-decoder baseline competitive on
    several ETT benchmarks \cite{das2023tide}.
    \item \textbf{TimeXer:} a Transformer-based model with cross-attention for
    exogenous variables, competitive on Weather and ETTh2 \cite{wang2024timexer}.
    \item \textbf{Pointwise+Refine (formerly labeled BERTOnly):} an ablation-style baseline that applies deeper encoder
    refinement without the patch formulation.
    \item \textbf{PatchTST$_\text{cap}$:} a capacity-matched PatchTST baseline scaled
    to approximately match \modelvzero\ parameter counts at each horizon, used to
    isolate the fusion contribution from model capacity (Section~\ref{sec:capmatch}).
\end{itemize}
The main comparison in Table~\ref{tab:main_results} includes DLinear, PatchTST,
iTransformer, TiDE, TimeXer, \modelvzero, and \modelvtwo. Pointwise+Refine is reported as
an ablation baseline (Table~\ref{tab:ablation}). PatchTST$_\text{cap}$ is reported
in the capacity-matched control analysis (Table~\ref{tab:capmatch}).

\subsection{Implementation and Training Protocol}
All models are implemented in PyTorch and trained under a shared 
evaluation pipeline so that split logic, horizon definitions, and metric 
computation remain consistent across comparisons.
All experiments use a unified training protocol:
\begin{itemize}
    \item Optimizer: Adam; \textbf{learning rate $10^{-4}$} for all
    patch-based models (PatchTST, PFBv0, PFBv2, PatchTST$_\text{cap}$);
    \textbf{learning rate $10^{-2}$} for DLinear following its canonical implementation
    \item Batch size: 32
    \item Epochs: 100 with early stopping (patience=10) on validation MSE
    \item Model selection: for each run, the checkpoint with best 
    validation MSE is evaluated on the test split
    \item Patch parameters (patch-based models): patch length $P=16$, 
    stride $S=8$
    \item Fusion/backbone dimensions: $d_{\text{model}}=128$, 
    $d_{\text{ff}}=512$, 8 attention heads, dropout 0.2 (applied after
    attention and feed-forward sublayers), no weight decay, no gradient clipping,
    no warmup schedule
\end{itemize}
iTransformer, TiDE, and TimeXer are evaluated under the same shared pipeline
with learning rate $10^{-3}$, consistent with their reported training settings.
Under the unified protocol, the capacity-matched control indicates that both
fusion variants have lower error than the matched PatchTST baseline on ETTh2
and Weather at every horizon, while on ETTm2 \modelvzero\ remains the
lowest-MSE fusion variant but \modelvtwo\ is not uniformly better at $H=336$.
The evidence against a pure capacity explanation is therefore most consistent on
Weather and ETTh2 and remains qualified, rather than universal, on ETTm2.
The reported \model\ variants differ only in the post-refinement fusion 
rule (Concat in \modelvzero\ vs Concat+Proj in \modelvtwo), allowing 
performance differences between them to be attributed primarily to the 
integration mechanism.

\subsection{Preprocessing and Normalization}
To reduce preprocessing confounds, all models use identical
split definitions, horizon settings, and train-split scaling
pipelines for each dataset. To ensure a fair comparison,
RevIN is treated as a fixed model-specific architectural
component rather than a tunable advantage~\cite{kim2022revin}:
it is enabled for \patchtst\ and both \model\ variants
following their original implementations, and disabled for
\dlinear\ and iTransformer following theirs. This respects
each model's intended design while maintaining a consistent
evaluation pipeline across all comparisons.

\subsection{Multi-Seed Evaluation and Variability}
We include representative multi-seed evaluation using five independent seeds
(2021--2025) for the $H=192$ stability analysis across the six standard datasets.
In addition, $H=96$ and $H=336$ multi-seed summaries are used to support
horizon-level trends, while capacity-control and Illness follow-up experiments use
three runs under the aligned LR=$10^{-4}$ protocol. We report mean, standard
deviation, and coefficient of variation where applicable, and interpret
per-dataset comparisons cautiously because five-seed per-configuration tests remain
underpowered for formal significance testing.

\subsection{Additional Analyses}
Beyond core benchmark comparisons, we include pooled statistical tests,
capacity-matched controls, K-depth sensitivity, component ablations, a controlled
efficiency benchmark, broader robustness probes, CKA representation diagnostics,
dataset-regime characterization, and patch-length sensitivity analysis.

\subsection{Reproducibility and Scope}
The protocol is designed for reproducibility through fixed splits, fixed horizons,
explicit optimization settings, and a shared evaluation pipeline. All reported table
values are synchronized to a single locked experiment snapshot.
\textbf{Hardware and software:} All experiments were run on an NVIDIA RTX 2080 Ti
(11\,GB VRAM) under CUDA 11.8, PyTorch 2.0.1, Python 3.10, Ubuntu 20.04.
No per-baseline canonical hyperparameter overrides were applied; the shared
protocol is the same across all compared models within each table.
All reported values come from this locked result archive. Random seeds control model
initialization and data-loader shuffling where applicable: the main single-seed
benchmark uses seed 2021; multi-seed stability analyses (Table~\ref{tab:multiseed}
and Appendix~\ref{app:full_results}) use five seeds (2021--2025); capacity-matched
control and Illness capacity-control follow-up experiments use seeds 2021, 2022,
and 2023. Scaling parameters are fit on the training split and then applied to
validation/test splits, while RevIN is used only where it is part of the original
model implementation. The code snapshot used for this manuscript corresponds to
commit e2b5fbb, and the repository includes the locked result files from which the
manuscript tables are assembled.

% =========================
\section{Results and Analysis}\label{sec:results}
% =========================

This section evaluates \model\ under the benchmark protocol and interprets outcomes
with emphasis on recurring patterns rather than isolated wins. We give greater weight
to consistent trends across datasets and horizons, relative improvements, and
representative multi-seed variability.

\subsection{Main Benchmark Performance}
Table~\ref{tab:main_results} summarizes benchmark performance. Overall, \model\ is
competitive, with lower errors in several regimes, but is not uniformly better across datasets.
Weather is the most favorable case. ETTm2 is also somewhat favorable, though the
capacity-matched comparison is mixed at $H=336$ for \modelvtwo. ETTh2 shows a similar
anti-capacity pattern within the PatchTST family, but TiDE has lower error overall at
longer horizons. ETTh1 and ETTm1 are neutral; several baselines, including DLinear,
remain competitive. Exchange is a known challenging benchmark for complex deep models
where fusion provides limited advantage. These structured outcomes support the paper's
regime-dependent interpretation. Figure~\ref{fig:perf_heatmap} complements the
single-seed main benchmark by showing five-seed relative MSE patterns for
\modelvzero\ against PatchTST and DLinear on the six standard non-Illness datasets.

\begin{figure*}[!t]
\centering
\includegraphics[width=\textwidth]{figures-300426/Fig2_mse_heatmap.png}
\caption{Relative MSE of \modelvzero\ compared with PatchTST and DLinear across
the six standard non-Illness datasets and three locked horizons, using five-seed
mean results. Negative values indicate lower MSE for \modelvzero; positive
values indicate lower MSE for the baseline. Illness is excluded because it uses
a different horizon set and a different error scale. TiDE and TimeXer are
reported in Table~\ref{tab:main_results} but omitted from this heatmap for
readability.}
\label{fig:perf_heatmap}
\end{figure*}

\begin{table*}[!t]
\centering
\caption{Forecasting performance across datasets and horizons (reported as MSE/MAE).
Bold indicates the lowest value overall, underline highlights runner-up or comparison-relevant values discussed in text.
\textbf{All values are from a single-seed run (seed 2021) under the unified
protocol (patch models LR=$10^{-4}$; DLinear LR=$10^{-2}$; iTransformer, TiDE,
TimeXer LR=$10^{-3}$); representative multi-seed variability (five seeds: 2021--2025)
is reported separately in Table~\ref{tab:multiseed}.}}
\label{tab:main_results}
\resizebox{\textwidth}{!}{%
\begin{tabular}{@{}llccc@{}}
\toprule
Dataset & Model & $H=96$ (MSE/MAE) & $H=192$ (MSE/MAE) & $H=336$ (MSE/MAE) \\
\midrule
\multirow{7}{*}{ETTh1}
& \dlinear      & 0.3711 / \textbf{0.3927} & \textbf{0.4043} / \textbf{0.4128} & \textbf{0.4345} / \textbf{0.4352} \\
& \patchtst     & 0.3814 / 0.4063 & 0.4312 / 0.4325 & 0.4934 / 0.4901 \\
& iTransformer  & 0.3997 / 0.4174 & 0.4495 / 0.4506 & 0.4636 / 0.4645 \\
& TiDE          & 0.3958 / 0.4141 & 0.4283 / 0.4327 & 0.4502 / 0.4466 \\
& TimeXer       & 0.3909 / 0.4076 & 0.4239 / 0.4269 & 0.4575 / 0.4532 \\
& \modelvzero   & \textbf{\underline{0.3675}} / \underline{0.3952} & \underline{0.4106} / \underline{0.4242} & \underline{0.4378} / 0.4476 \\
& \modelvtwo    & 0.3863 / 0.4053 & 0.4198 / 0.4266 & 0.4411 / \underline{0.4404} \\
\midrule
\multirow{7}{*}{ETTh2}
& \dlinear      & \textbf{0.2836} / \textbf{0.3473} & 0.3797 / 0.4151 & 0.4257 / 0.4555 \\
& \patchtst     & 0.3145 / 0.3647 & 0.3887 / 0.4118 & 0.4059 / 0.4293 \\
& iTransformer  & 0.3062 / 0.3606 & 0.3667 / 0.4009 & 0.3988 / 0.4229 \\
& TiDE          & 0.2914 / 0.3515 & \textbf{0.3517} / \textbf{0.3903} & \textbf{0.3736} / \textbf{0.4109} \\
& TimeXer       & 0.2969 / 0.3604 & 0.3614 / 0.3972 & 0.3945 / 0.4265 \\
& \modelvzero   & 0.3054 / 0.3524 & 0.3720 / 0.4004 & \underline{0.3814} / \underline{0.4139} \\
& \modelvtwo    & \underline{0.2893} / \underline{0.3474} & \underline{0.3673} / \underline{0.3952} & 0.3955 / 0.4225 \\
\midrule
\multirow{7}{*}{ETTm1}
& \dlinear      & 0.3003 / \textbf{0.3448} & 0.3357 / \textbf{0.3663} & 0.3705 / \textbf{0.3872} \\
& \patchtst     & \textbf{0.2982} / 0.3490 & 0.3429 / 0.3776 & 0.3741 / 0.3981 \\
& iTransformer  & 0.3099 / 0.3613 & 0.3424 / 0.3801 & 0.3786 / 0.4015 \\
& TiDE          & 0.3069 / 0.3485 & 0.3416 / 0.3690 & 0.3811 / 0.3929 \\
& TimeXer       & 0.3160 / 0.3614 & 0.3612 / 0.3862 & 0.3863 / 0.4052 \\
& \modelvzero   & 0.3037 / 0.3479 & 0.3426 / 0.3747 & 0.3823 / 0.3991 \\
& \modelvtwo    & 0.2984 / 0.3466 & \textbf{0.3326} / \underline{0.3690} & \textbf{0.3695} / 0.3961 \\
\midrule
\multirow{7}{*}{ETTm2}
& \dlinear      & \textbf{0.1653} / 0.2574 & \underline{0.2273} / 0.3077 & 0.2930 / 0.3524 \\
& \patchtst     & 0.1728 / 0.2640 & 0.2282 / 0.3018 & 0.2935 / 0.3457 \\
& iTransformer  & 0.1749 / 0.2660 & 0.2475 / 0.3139 & 0.3020 / 0.3484 \\
& TiDE          & 0.1690 / 0.2595 & \textbf{0.2241} / \textbf{0.2965} & 0.2780 / 0.3314 \\
& TimeXer       & 0.1716 / \textbf{0.2562} & 0.2311 / 0.2986 & 0.2870 / 0.3351 \\
& \modelvzero   & \underline{0.1679} / 0.2571 & 0.2347 / \underline{0.2997} & \textbf{0.2739} / \textbf{0.3291} \\
& \modelvtwo    & \underline{0.1695} / \underline{0.2567} & 0.2347 / 0.3001 & 0.3014 / 0.3496 \\
\midrule
\multirow{7}{*}{Exchange}
& \dlinear      & \textbf{0.0849} / \textbf{0.2052} & \textbf{0.1539} / \textbf{0.2860} & \textbf{0.2890} / \textbf{0.4027} \\
& \patchtst     & 0.0953 / \underline{0.2195} & 0.1913 / 0.3182 & \underline{0.3139} / \underline{0.4092} \\
& iTransformer  & \underline{0.0949} / 0.2210 & 0.2138 / 0.3358 & 0.3750 / 0.4523 \\
& TiDE          & 0.0987 / 0.2228 & 0.1925 / 0.3143 & 0.3422 / 0.4261 \\
& TimeXer       & 0.1050 / 0.2259 & 0.1911 / 0.3124 & 0.3792 / 0.4486 \\
& \modelvzero   & 0.0997 / 0.2224 & 0.2184 / 0.3340 & 0.3533 / 0.4340 \\
& \modelvtwo    & 0.0991 / 0.2209 & \underline{0.1878} / \underline{0.3108} & 0.3381 / 0.4223 \\
\midrule
\multirow{7}{*}{Weather}
& \dlinear      & 0.1782 / 0.2432 & 0.2179 / 0.2775 & 0.2612 / 0.3116 \\
& \patchtst     & 0.1526 / 0.2019 & 0.1962 / 0.2429 & 0.2479 / 0.2814 \\
& iTransformer  & 0.1603 / 0.2096 & 0.2031 / 0.2487 & 0.2527 / 0.2867 \\
& TiDE          & 0.1764 / 0.2267 & 0.2185 / 0.2612 & 0.2659 / 0.2957 \\
& TimeXer       & 0.1508 / 0.2024 & \underline{0.1942} / \underline{0.2409} & \underline{0.2453} / 0.2811 \\
& \modelvzero   & \textbf{0.1496} / \underline{0.2007} & \textbf{0.1936} / \textbf{0.2387} & 0.2517 / 0.2839 \\
& \modelvtwo    & \underline{0.1522} / \textbf{0.2000} & 0.1938 / 0.2394 & \textbf{0.2446} / \textbf{0.2802} \\
\bottomrule
\end{tabular}%
}
\end{table*}

Aggregate trend views are shown in Figure~\ref{fig:dataset_performance} (dataset-wise
means) and Figure~\ref{fig:horizon_scaling} (horizon-wise means).

\begin{figure*}[!t]
\centering
\includegraphics[width=\textwidth]{figures/fig3_dataset_performance.png}
\caption{Dataset-level mean MSE across horizons ($H \in \{96,192,336\}$) for the
main comparison models. Illness uses different horizons ($H=24,48,60$) and is
excluded here for scale comparability. Lower is better.}
\label{fig:dataset_performance}
\end{figure*}

\begin{figure*}[!t]
\centering
\includegraphics[width=\textwidth]{figures/fig4_horizon_scaling.png}
\caption{Horizon-scaling analysis showing mean MSE across the six standard
non-Illness datasets at each horizon ($H \in \{96,192,336\}$). Lower is better.}
\label{fig:horizon_scaling}
\end{figure*}

\subsection{Performance on ETTh1 and ETTh2}
On ETTh1, \model\ remains competitive, with \modelvzero\ achieving the lowest MSE
among the compared models at $H=96$. Under the corrected unified protocol
(DLinear LR=$10^{-2}$), DLinear achieves the lowest MSE at $H=192$ and $H=336$,
with \modelvzero\ as runner-up at $H=192$. Overall margins are small and rankings
vary with horizon.

ETTh2 shows a somewhat clearer pattern for fusion within the PatchTST family:
\modelvtwo\ ranks second at $H=96$ while DLinear now achieves the lowest MSE there
(0.2836). \modelvzero\ ranks second at $H=336$. However,
TiDE has the lowest MSE on ETTh2 at $H=192$ and $H=336$, and this is
an important qualification: ETTh2 supports the anti-capacity argument within the
PatchTST family (Section~\ref{sec:capmatch}), but does not support a claim of
uniformly lower error across all architectures.

\subsection{Performance on ETTm1 and ETTm2}
ETTm1 is neutral. Multiple models DLinear, PatchTST, and \modelvtwo achieve
comparable performance at $H=96$ and $H=336$. The gains from fusion are limited, and
DLinear has the lowest or tied-lowest MSE at the lowest horizon.

ETTm2 is somewhat favorable but horizon-dependent. In the main benchmark,
DLinear achieves the lowest MSE at $H=96$ (0.1653), while \modelvzero\
achieves the lowest MSE at $H=336$ (0.2739); TiDE, PatchTST, and TimeXer have
lower MSE than the fusion variants at $H=192$. In the five-seed
unified-protocol capacity-control study, however, both fusion variants outperform
PatchTST$_\text{base}$ and PatchTST$_\text{cap}$ at $H=192$, and \modelvzero\
remains the lowest-MSE fusion variant at $H=336$. The main qualification is
therefore specific to the main-benchmark $H=192$ setting and the $H=336$
comparison for \modelvtwo.

\subsection{Performance on Exchange}
Exchange is a challenging benchmark for complex deep models and serves as a useful
negative control for the fusion architecture. Under the corrected DLinear protocol
(LR=$10^{-2}$), DLinear achieves the lowest MSE at all three horizons, with
\modelvzero\ consistently worse than the best non-fusion baseline. The
representative $H=192$ multi-seed results further support this interpretation. This result is still
informative: it bounds the scope of the contribution and supports a measured rather
than over-claimed interpretation.

\subsection{Performance on Weather}
Weather is the most favorable dataset overall. \modelvtwo\ is relatively favorable,
including the lowest main-table MSE at $H=336$, while \modelvzero\ is lowest at
the shorter horizons and remains competitive, though not top-two, at $H=336$ in the
main benchmark. The multi-seed table shows that \modelvzero\ (mean $0.1964$) and
\modelvtwo\ (mean $0.1941$) both have slightly lower mean MSE than PatchTST (mean $0.1976$) at
$H=192$. Under the aligned capacity-control follow-up, both fusion variants also
have lower error than PatchTST$_\text{base}$ and PatchTST$_\text{cap}$ across all three
horizons. Supplementary diagnostic views help localize this advantage. In
Appendix~\ref{app:full_results}, Figure~\ref{fig:app_error_by_step} shows that on
Weather the advantage over DLinear becomes clearer at medium-to-long lead times,
while Figure~\ref{fig:app_per_variable_weather} shows that the gains relative to
PatchTST are concentrated on smoother thermodynamic channels rather than the sparse
precipitation-related channels.

\subsection{Representative Multi-Seed Variability}
Table~\ref{tab:multiseed} reports mean$\pm$std across five seeds (2021--2025) for representative
$H=192$ settings, with coefficient of variation (CV)
computed on MSE. The proposed variants do not show catastrophic
instability, but seed sensitivity is dataset-dependent.
Accordingly, we avoid broad statistical-dominance claims over the full
21-setting grid. The multi-seed table is used to contextualize uncertainty
in representative regimes and to avoid over-interpreting very small gaps
in the main benchmark table. Notably, the CV for \modelvzero\ remains
below 2\% on both ETTm2 (CV=1.84\%) and Weather (CV=0.97\%), indicating
that the reported gains on those datasets are consistent across the tested
initializations and less likely to be single-run artifacts. A complementary 95\%
confidence-interval view is given in Figure~\ref{fig:ci_errorbars}. It visually
reinforces the descriptive reading: separation from DLinear is clearest on Weather
and ETTh2, whereas ETTm2 remains favorable but closer. For qualitative
trajectory-level seed dispersion rather than calibrated predictive uncertainty,
Appendix~\ref{app:full_results} also provides illustrative forecast-dispersion
plots in Figure~\ref{fig:app_forecast_dispersion}.
Per-dataset five-seed comparisons are treated descriptively because they are
underpowered for formal per-configuration significance testing. Formal hypothesis
tests are therefore reported only at the pooled and horizon-stratified levels in
Section~\ref{sec:wilcoxon}, and are interpreted as global supporting evidence
rather than proof of superiority on any individual dataset.

\begin{figure*}[!t]
\centering
\includegraphics[width=\textwidth]{figures-300426/Fig_CI_errorbars.png}
\caption{Five-seed MSE means with 95\% confidence intervals for the representative
$H=192$ standard-dataset settings. Intervals are computed from the same five-seed
runs summarized in Table~\ref{tab:multiseed} using the $t$-distribution with four
degrees of freedom. Lower is better. The clearest separation from DLinear appears
on Weather and ETTh2, whereas ETTm2 remains favorable but closer.}
\label{fig:ci_errorbars}
\end{figure*}

\subsection{Pooled Statistical Significance Testing}
\label{sec:wilcoxon}
Per-configuration tests with $N=5$ seeds have minimum achievable Wilcoxon
$p=0.0625$, so no individual 5-seed comparison can reach $\alpha=0.05$
significance. We therefore pool paired observations across nine 5-seed
configurations ($N=45$ pairs: six standard datasets at $H=192$ plus Illness
$H\in\{24,48,60\}$) and apply two-sided Wilcoxon signed-rank tests with
Holm--Bonferroni and Benjamini--Hochberg FDR correction (8 simultaneous
tests).

\textbf{Key results.} Results are classified into three tiers: (1)~Holm-corrected
significant, (2)~BH-FDR significant but not Holm-corrected, (3)~directional but not
corrected-significant.

PFBv0 vs.\ DLinear: \textit{Holm-significant} for both MSE
($p_\text{holm}=0.027$) and MAE ($p_\text{holm}=0.00039$), with
rank-biserial effect sizes $r=-0.49$ (MSE) and $r=-0.66$ (MAE) ---
moderate-to-large. PFBv0 vs.\ PatchTST: MSE is \textit{directional only}
($p_\text{holm}=0.17$; raw $p=0.035$); MAE is \textit{BH-FDR significant}
($p_\text{fdr}=0.031$) but does not survive Holm correction.

A horizon-stratified analysis over $N=30$ pairs (six standard datasets
$\times$ 5 seeds) adds further nuance: at $H=96$, PFBv0 and PFBv2
\textit{Holm-significantly} outperform PatchTST on MAE; at $H=336$,
PFBv0 \textit{Holm-significantly} outperforms both PatchTST (MSE and MAE)
and DLinear (MAE); DLinear MSE at $H=336$ is \textit{BH-FDR significant}
but not Holm-corrected.

A compact visual summary of the pooled and horizon-stratified Wilcoxon outcomes is
provided in Appendix~\ref{app:full_results}, Figure~\ref{fig:app_significance}.

\textbf{Interpretation caveats.} These pooled tests use Illness horizons
alongside standard horizons, which may increase power artificially if Illness
has very different scale. The pooled result is best treated as global
supporting evidence that the directional pattern is not a systematic
single-run artefact, rather than as proof of superiority on any specific
dataset. Dataset-level conclusions continue to rely on directional tables,
multi-seed variability, and capacity-matched controls.

\begin{table}[t]
\centering
\caption{Representative multi-seed validation (mean $\pm$ std) for $H=192$
forecasts under the unified protocol (five seeds: 2021--2025; DLinear
LR=$10^{-2}$, patch models LR=$10^{-4}$).
CV = coefficient of variation (\%) on MSE.}
\label{tab:multiseed}
\small
\resizebox{\columnwidth}{!}{%
\begin{tabular}{@{}llccc@{}}
\toprule
Dataset & Model & MSE & MAE & CV (\%) \\
\midrule
\multirow{4}{*}{ETTm1}
& DLinear    & \textbf{0.3351 $\pm$ 0.0020} & \textbf{0.3658 $\pm$ 0.0032} & 0.60 \\
& PatchTST   & 0.3364 $\pm$ 0.0041 & 0.3753 $\pm$ 0.0020 & 1.22 \\
& \modelvzero & 0.3395 $\pm$ 0.0023 & 0.3743 $\pm$ 0.0018 & 0.68 \\
& \modelvtwo  & 0.3403 $\pm$ 0.0072 & 0.3753 $\pm$ 0.0056 & 2.12 \\
\midrule
\multirow{4}{*}{ETTm2}
& DLinear    & 0.2308 $\pm$ 0.0051 & 0.3104 $\pm$ 0.0065 & 2.21 \\
& PatchTST   & 0.2370 $\pm$ 0.0055 & 0.3070 $\pm$ 0.0037 & 2.32 \\
& \modelvzero & \textbf{0.2287 $\pm$ 0.0042} & \textbf{0.2975 $\pm$ 0.0023} & 1.84 \\
& \modelvtwo  & 0.2323 $\pm$ 0.0031 & 0.2986 $\pm$ 0.0024 & 1.33 \\
\midrule
\multirow{4}{*}{ETTh1}
& DLinear    & \textbf{0.4074 $\pm$ 0.0023} & \textbf{0.4182 $\pm$ 0.0029} & 0.56 \\
& PatchTST   & 0.4321 $\pm$ 0.0014 & 0.4360 $\pm$ 0.0047 & 0.32 \\
& \modelvzero & 0.4196 $\pm$ 0.0101 & 0.4297 $\pm$ 0.0065 & 2.41 \\
& \modelvtwo  & 0.4225 $\pm$ 0.0071 & 0.4308 $\pm$ 0.0066 & 1.68 \\
\midrule
\multirow{4}{*}{ETTh2}
& DLinear    & 0.3945 $\pm$ 0.0072 & 0.4259 $\pm$ 0.0052 & 1.82 \\
& PatchTST   & 0.3824 $\pm$ 0.0071 & 0.4084 $\pm$ 0.0030 & 1.86 \\
& \modelvzero & 0.3740 $\pm$ 0.0083 & 0.3980 $\pm$ 0.0028 & 2.22 \\
& \modelvtwo  & \textbf{0.3699 $\pm$ 0.0065} & \textbf{0.3963 $\pm$ 0.0032} & 1.76 \\
\midrule
\multirow{4}{*}{Exchange}
& DLinear    & \textbf{0.1779 $\pm$ 0.0056} & \textbf{0.3160 $\pm$ 0.0046} & 3.15 \\
& PatchTST   & 0.1921 $\pm$ 0.0052 & 0.3180 $\pm$ 0.0050 & 2.71 \\
& \modelvzero & 0.2244 $\pm$ 0.0084 & 0.3395 $\pm$ 0.0087 & 3.75 \\
& \modelvtwo  & 0.2011 $\pm$ 0.0133 & 0.3217 $\pm$ 0.0109 & 6.61 \\
\midrule
\multirow{4}{*}{Weather}
& DLinear    & 0.2168 $\pm$ 0.0020 & 0.2757 $\pm$ 0.0042 & 0.92 \\
& PatchTST   & 0.1976 $\pm$ 0.0016 & 0.2440 $\pm$ 0.0017 & 0.81 \\
& \modelvzero & 0.1964 $\pm$ 0.0019 & 0.2429 $\pm$ 0.0022 & 0.97 \\
& \modelvtwo  & \textbf{0.1941 $\pm$ 0.0014} & \textbf{0.2409 $\pm$ 0.0014} & 0.72 \\
\bottomrule
\end{tabular}
}
\end{table}

\subsection{Illness as a Low-Data Case Study}
\label{sec:illness}
Table~\ref{tab:illness_standard} reports the standard single-seed benchmark on
Illness. Under that protocol, \modelvzero\ is the lowest-MSE model across the standard
baseline set at $H=24$, $H=48$, and $H=60$. This supports only a narrow claim: against the
standard single-seed baselines, \modelvzero\ has the lowest MSE on Illness.

Table~\ref{tab:illness_capacity} reports the aligned 3-seed capacity-control
follow-up as mean$\pm$std (LR=$10^{-4}$). This follow-up changes the interpretation.
PatchTST$_\text{cap}$ has the lowest mean MSE at $H=24$ and $H=48$, while
PatchTST$_\text{base}$ has the lowest mean MSE at $H=60$. The more cautious interpretation is therefore
that Illness supports the value of additional patch-based capacity in a low-data
regime, but it does not cleanly isolate a fusion-specific advantage.

\begin{table}[t]
\centering
\caption{Illness standard single-seed benchmark (unified protocol: DLinear
LR=$10^{-2}$, patch models LR=$10^{-4}$). Lower MSE is better.}
\label{tab:illness_standard}
\small
\resizebox{\columnwidth}{!}{%
\begin{tabular}{@{}lccc@{}}
\toprule
Model & $H=24$ & $H=48$ & $H=60$ \\
\midrule
\dlinear               & 2.6181 & 2.2883 & 2.4886 \\
\patchtst              & 2.3322 & 1.9039 & 2.1090 \\
iTransformer           & 2.2531 & 2.2221 & 2.4974 \\
TiDE                   & 3.2519 & 3.5216 & 3.8371 \\
TimeXer                & 2.2997 & 2.1596 & 2.8440 \\
\modelvzero            & \textbf{1.9868} & \textbf{1.8855} & \textbf{1.9930} \\
\modelvtwo             & 2.2976 & 2.0007 & 2.9480 \\
Pointwise+Refine       & 3.6351 & 3.7898 & 4.7495 \\
\bottomrule
\end{tabular}
}
\end{table}

\begin{table}[t]
\centering
\caption{Illness capacity-control follow-up (MSE mean $\pm$ std over 3 seeds,
LR=$10^{-4}$). To avoid protocol mixing, only models run under the aligned
follow-up setting are shown. Lower is better.}
\label{tab:illness_capacity}
\small
\resizebox{\columnwidth}{!}{%
\begin{tabular}{@{}lccc@{}}
\toprule
Model & $H=24$ & $H=48$ & $H=60$ \\
\midrule
PatchTST$_\text{base}$ & 2.076 $\pm$ 0.132 & 2.078 $\pm$ 0.154 & \textbf{2.000 $\pm$ 0.019} \\
PatchTST$_\text{cap}$  & \textbf{1.899 $\pm$ 0.021} & \textbf{1.892 $\pm$ 0.003} & 2.097 $\pm$ 0.099 \\
\modelvzero            & 2.061 $\pm$ 0.169 & \underline{1.911 $\pm$ 0.024} & 2.036 $\pm$ 0.071 \\
\modelvtwo             & 2.200 $\pm$ 0.148 & 2.058 $\pm$ 0.061 & 2.072 $\pm$ 0.016 \\
\bottomrule
\end{tabular}
}
\end{table}

\subsection{Parameter-Matched Capacity Control}
\label{sec:capmatch}
A key question in interpreting any performance gain from a larger model is whether
the gain comes from the proposed architectural design or simply from having more
parameters. To address this directly, we trained a capacity-matched PatchTST variant
(PatchTST$_\text{cap}$) by increasing model width until its parameter count
approximately matched \modelvzero\ at each horizon. All other hyperparameters were
held constant.

Table~\ref{tab:capmatch} shows results on ETTh2, ETTm2, and Weather, the three
datasets outside of Illness where the evidence for fusion is most favorable. A recurring pattern is
that PatchTST$_\text{cap}$ does not fully close the gap and in some cases performs
worse than PatchTST$_\text{base}$. On ETTh2 and Weather, both fusion variants
have lower error than the capacity-matched baseline at every horizon. On ETTm2, both variants
have lower error at $H=96$ and $H=192$, while at $H=336$ \modelvzero\ remains lowest and
PatchTST$_\text{cap}$ is slightly lower than \modelvtwo. This is most consistent with the view
that the anti-capacity evidence is most consistent on Weather and ETTh2 and qualified,
rather than universal, on ETTm2. Figure~\ref{fig:capmatch}
shows the corresponding compact $H=192$ view.

The Illness result is more complex and is discussed in Section~\ref{sec:illness}.

\begin{figure}[t]
\centering
\includegraphics[width=0.98\linewidth]{figures-300426/Fig3_capacity_control.png}
\caption{Capacity-matched control results on selected $H=192$ settings under the
3-run median protocol (LR=$10^{-4}$). Scaling PatchTST to matched
parameter counts does not reproduce the fusion gains in these selected cases,
which is consistent with an effect not explained by width scaling alone.}
\label{fig:capmatch}
\end{figure}

\begin{table*}[t]
\centering
\caption{Capacity-matched control experiment (median over 3 runs, follow-up
protocol LR=$10^{-4}$). PatchTST$_\text{cap}$ is scaled to approximately match
\modelvzero\ parameter counts. PatchTST$_\text{depth}$ is a depth-matched variant
(\textasciitilde3.22M params at $H=192$, deeper stacking) evaluated at $H=192$ only
(---\ indicates horizon not tested in this control). The Params column reports exact,
horizon-specific follow-up protocol counts in millions and should not be compared one-to-one with the representative parameter
values in Table~\ref{tab:architecture}. Values are not numerically identical to
the main benchmark protocol. Bold indicates the lowest value per horizon, underline indicates
second-lowest per horizon.}
\label{tab:capmatch}
\resizebox{\textwidth}{!}{%
\begin{tabular}{@{}llccccc@{}}
\toprule
Dataset & Model & Params ($H=96/192/336$, M) & $H=96$ MSE & $H=192$ MSE & $H=336$ MSE \\
\midrule
\multirow{5}{*}{ETTh2}
& PatchTST$_\text{base}$ & 1.11 / 1.63 / 2.40 & 0.3179 & 0.3848 & 0.4075 \\
& PatchTST$_\text{cap}$  & 2.21 / 3.28 / 4.81 & 0.3155 & 0.3883 & 0.4238 \\
& PatchTST$_\text{depth}$ & --- / 3.22 / --- & --- & 0.3833 & --- \\
& \modelvzero             & 2.22 / 3.26 / 4.81 & \underline{0.3007} & \underline{0.3709} & \textbf{0.3868} \\
& \modelvtwo              & 1.94 / 2.46 / 3.23 & \textbf{0.2972} & \textbf{0.3701} & \underline{0.3945} \\
\midrule
\multirow{5}{*}{ETTm2}
& PatchTST$_\text{base}$ & 1.11 / 1.63 / 2.40 & 0.1770 & 0.2348 & 0.2950 \\
& PatchTST$_\text{cap}$  & 2.21 / 3.28 / 4.81 & 0.1811 & 0.2374 & \underline{0.2898} \\
& PatchTST$_\text{depth}$ & --- / 3.22 / --- & --- & 0.2413 & --- \\
& \modelvzero             & 2.22 / 3.26 / 4.81 & \textbf{0.1707} & \textbf{0.2287} & \textbf{0.2821} \\
& \modelvtwo              & 1.94 / 2.46 / 3.23 & \underline{0.1717} & \underline{0.2318} & 0.2979 \\
\midrule
\multirow{5}{*}{Weather}
& PatchTST$_\text{base}$ & 1.11 / 1.63 / 2.40 & 0.1515 & 0.1977 & 0.2539 \\
& PatchTST$_\text{cap}$  & 2.21 / 3.28 / 4.81 & 0.1546 & 0.2026 & 0.2565 \\
& PatchTST$_\text{depth}$ & --- / 3.22 / --- & --- & 0.1969 & --- \\
& \modelvzero             & 2.22 / 3.26 / 4.81 & \textbf{0.1492} & \underline{0.1957} & \underline{0.2492} \\
& \modelvtwo              & 1.94 / 2.46 / 3.23 & \underline{0.1496} & \textbf{0.1946} & \textbf{0.2481} \\
\bottomrule
\end{tabular}
}
\end{table*}

\subsection{K-Depth Ablation}
\label{sec:kdepth}
The post-encoder refinement stack depth $K$ was fixed at $3$ in the main benchmark
experiments. To assess whether this choice is justified, we trained \modelvzero\ with
$K \in \{1,2,3\}$ on ETTh2 and Weather across all three horizons using a single fixed
seed. Results are shown in Figure~\ref{fig:kdepth} and Table~\ref{tab:kdepth}.

The main finding is that differences between $K$ values are small and inconsistent.
Across the six tested settings, $K=1$ achieves the lowest MSE in four cases and
$K=2$ in two. Because these are single-seed follow-up runs, the ablation is
interpreted as descriptive MSE evidence rather than a formal significance test.

The cautious conclusion is that the existence of a refinement stage matters more than
its depth within this range. \textbf{$K=1$ is consistent with capturing most of
the observed benefit in the tested single-seed setting}, and offers lower
computational cost. We therefore treat this as descriptive evidence that $K=1$ may
be adequate as a default, while noting that multi-seed replication on further datasets
would be needed to strengthen this recommendation. $K=3$ was
retained in main experiments for consistency with the original design.

\begin{figure}[t]
\centering
\includegraphics[width=0.98\linewidth]{figures-300426/A5_kdepth_sensitivity.png}
\caption{K-depth ablation: MSE versus refinement depth $K \in \{1,2,3\}$ on ETTh2
and Weather under the follow-up protocol (single seed, LR=$10^{-4}$).
Performance is largely insensitive to $K$ within this range, and $K=1$
captures most of the architectural benefit.}
\label{fig:kdepth}
\end{figure}

\begin{table}[t]
\centering
\caption{K-depth ablation for \modelvzero\ (single seed, LR=$10^{-4}$).
Lower MSE is better. Bold indicates the lowest value per setting.}
\label{tab:kdepth}
\small
\resizebox{\columnwidth}{!}{%
\begin{tabular}{@{}llccc@{}}
\toprule
Dataset & Horizon & $K=1$ MSE & $K=2$ MSE & $K=3$ MSE \\
\midrule
\multirow{3}{*}{ETTh2}
& $H=96$  & \textbf{0.3081} & 0.3100 & 0.3092 \\
& $H=192$ & \textbf{0.3727} & 0.3782 & 0.3733 \\
& $H=336$ & 0.3967 & \textbf{0.3892} & 0.3931 \\
\midrule
\multirow{3}{*}{Weather}
& $H=96$  & \textbf{0.1493} & 0.1503 & 0.1494 \\
& $H=192$ & 0.1958 & \textbf{0.1942} & 0.1978 \\
& $H=336$ & \textbf{0.2495} & 0.2518 & 0.2506 \\
\bottomrule
\end{tabular}
}
\end{table}

\subsection{Interpreting v0 versus v2}
The comparison between \modelvzero\ (Concat) and \modelvtwo\ (Concat+Proj) suggests that
fusion design matters in addition to refinement depth. Neither variant is consistently
lower across all settings. On ETTm2, \modelvzero\ has lower MSE at $H=336$ while both variants tie at
$H=192$. On Weather, \modelvtwo\ has lower MSE at $H=336$. On ETTh2, \modelvtwo\ has lower MSE at
$H=96$ while \modelvzero\ has lower MSE at $H=336$. This pattern suggests treating fusion
design as a meaningful architectural component rather than only an implementation detail.

\subsection{Partial Ablation Evidence}
Table~\ref{tab:ablation} provides directional evidence on the role of patching versus
fusion. The most decisive aggregate result is that the \emph{Pointwise+Refine} ablation
(applying the 3-layer refinement stack to non-patched, point-wise input tokens, without
explicit fusion; formerly labeled \bertonly) is not consistently
competitive: adding deeper encoder refinement without patch tokenization yields much
weaker aggregate performance (aggregate MSE 0.906) than the patch-based variants.
This does not imply that refinement without patching is never useful --- the Pointwise+Refine ablation is
occasionally competitive in isolated settings but the full benchmark suggests that
refinement becomes reliable only when coupled with patch tokenization and explicit
fusion. Adding fusion on top of patching then improves the aggregate summary, with
\modelvtwo\ achieving the best aggregate value (0.562).

A new head-capacity control, PatchTST$_\text{LargeHead}$ (3.45M params, trained with
identical hyperparameters at seed=2021), adds a two-layer MLP output head to the
standard PatchTST backbone, exceeding \modelvzero's total parameter count by 6\%.
Despite this additional output-side capacity, it achieves aggregate MSE 0.604 ---
worse than \modelvzero\ (0.599) --- and underperforms PFBv0 on all three evaluated
datasets: ETTm2 H=192 (0.2351 vs 0.2216), ETTh2 H=192 (0.3792 vs 0.3733), and
Weather H=192 (0.1989 vs 0.1978). This supports the interpretation that the
observed improvements are not explained by output-head capacity alone in this
limited diagnostic setting (single-seed, H=192, three datasets). Figure~\ref{fig:ablation_study}
provides a compact visual summary of the $H=192$ component-level ablation on
three representative datasets. The two aggregates
are not numerically comparable: Table~\ref{tab:ablation} summarizes a broader,
scale-dependent ablation archive, whereas Figure~\ref{fig:ablation_study} is used
as a descriptive visual aid for the same model family.

\begin{table}[t]
\centering
\caption{Ablation summary: Aggregate MSE values come from a broader,
scale-dependent ablation archive and are used only as directional evidence;
primary conclusions rely on per-dataset tables and these values should not be
compared directly with Figure~\ref{fig:ablation_study}.
The \emph{Pointwise+Refine} row (formerly labeled \emph{BERTOnly}) is a
\emph{patch-removal ablation}: it applies the 3-layer refinement stack to
pointwise (non-patched) input tokens and is not intended as a competitive
baseline, only as a directional check on the role of patch tokenization.
PatchTST$_\text{LargeHead}$ is a head-capacity control (3.45M params,
trained with $d_{\text{ff}}=512$, LR=$10^{-4}$, seed=2021) that adds a
2-layer MLP output head to the standard PatchTST backbone; it exceeds
\modelvzero's 3.26M parameter count and serves as diagnostic evidence that
a larger output head alone does not reproduce the observed gains in the
tested $H=192$ settings.}
\label{tab:ablation}
\small
\resizebox{\columnwidth}{!}{%
\begin{tabular}{@{}lcccc@{}}
\toprule
Architecture & Patching & Refinement / Fusion & Aggregate MSE \\
\midrule
\dlinear    & \xmark            & \xmark & 0.842 \\
Pointwise+Refine & \xmark       & \shortstack[l]{Refine only,\\no fusion} & 0.906 \\
\patchtst   & \cmark\ ($P$=16) & \xmark & 0.622 \\
PatchTST$_\text{LargeHead}$ & \cmark\ ($P$=16) & \shortstack[l]{Large MLP head,\\no fusion} & 0.604 \\
\modelvzero & \cmark\ ($P$=16) & \shortstack[l]{Refine +\\Concat} & 0.599 \\
\modelvtwo  & \cmark\ ($P$=16) & \shortstack[l]{Refine +\\Concat+Proj} & \textbf{0.562} \\
\bottomrule
\end{tabular}
}
\end{table}

\begin{figure*}[!t]
\centering
\includegraphics[width=\textwidth]{figures-300426/Fig5_ablation_chain.png}
\caption{Single-seed component ablation at $H=192$ on ETTm2, ETTh2, and
Weather. The bars compare PatchOnly, RefineOnly, FusionOnly, and the full
\modelvzero\ configuration. Lower MSE is better. These results are diagnostic
and should be interpreted together with the broader aggregate ablation summary
in Table~\ref{tab:ablation}.}
\label{fig:ablation_study}
\end{figure*}

\subsection{Missing-Input Robustness Case Study}
We report a robustness case study by randomly zero-masking input entries at rates of
10\%, 20\%, and 30\% during inference on ETTm1 $H=192$ (no imputation module).
Masking is applied only at inference, independently across time steps and variables,
using the same mask realizations for all compared models; no imputation module,
retraining, or fine-tuning is used.
All three models degrade approximately linearly, but attention-based models degrade
more slowly than DLinear: DLinear reaches +27.4\% MSE degradation at 30\% missing,
while \modelvzero\ and PatchTST reach +22.3\% and +22.0\% respectively. The
difference between \modelvzero\ and PatchTST is negligible (0.17\% absolute gap at
$\rho=0.3$). Figure~\ref{fig:robustness} visualizes representative
random-missingness behavior on ETTm2 $H=192$ and Exchange $H=192$,
illustrating that robustness patterns are regime-dependent rather than uniformly
favorable to \model. This
analysis provides initial evidence that the fusion mechanism does not introduce
additional fragility under moderate input degradation, but does not establish broad
robustness across all datasets or missingness types.

To provide broader robustness context, we extended the stress-test to additional
corruption types and datasets: random missingness, block missingness, Gaussian
additive noise, and channel dropout on ETTm2, Weather, Exchange, and Illness.
All perturbations are applied at inference only with no retraining. Table~\ref{tab:robustness_summary}
summarizes the qualitative regime-level findings. Supplementary Gaussian-noise and
channel-dropout curves are shown in Appendix~\ref{app:full_results},
Figure~\ref{fig:app_gaussian_channel_robustness}. They reinforce that DLinear is
especially noise-robust on Weather, whereas channel dropout is severe for all
models on ETTm2.

\begin{table}[t]
\centering
\caption{Broader robustness regime summary (test-time perturbations, no retraining).
Each row reflects inference-only stress tests across the listed datasets. ``Best'' and
``leads'' are descriptive; no formal test is conducted.}
\label{tab:robustness_summary}
\small
\resizebox{\columnwidth}{!}{%
\begin{tabular}{@{}p{0.22\columnwidth}p{0.22\columnwidth}p{0.46\columnwidth}@{}}
\toprule
Corruption & Datasets & Qualitative finding \\
\midrule
Random missingness & ETTm2, Weather, Exchange, Illness &
PFBv0/PFBv2 competitive or best on ETTm2 $H=192$; DLinear strongest on
Exchange at all rates. \\
Block missingness & ETTm2, Exchange &
PFBv0$\approx$PFBv2 best on ETTm2 $H=96$; DLinear leads at $H=192$
and Exchange. \\
Gaussian noise & ETTm2, Weather &
DLinear highly robust on Weather; PFB shows no consistent noise
advantage across datasets. \\
Channel dropout & ETTm2, Weather &
No PFB robustness advantage; channel dropout severe for all models. \\
\bottomrule
\end{tabular}
}
\end{table}

\begin{figure}[!t]
\centering
\includegraphics[width=0.98\linewidth]{figures-300426/Fig6_robustness_curves.png}
\caption{Representative random-missingness stress tests on ETTm2 $H=192$ and
Exchange $H=192$. Relative MSE degradation at 10\%, 20\%, and 30\% random
zero-masking rates (inference only, no retraining). Lower is better. These
plots illustrate regime-dependent robustness behavior and are not formal
robustness superiority tests.}
\label{fig:robustness}
\end{figure}

\subsection{Patch-Length Sensitivity}
A sensitivity sweep over $P \in \{8,16,32,64\}$ on ETTm1 and ETTh1 for \modelvzero\
shows low sensitivity: MSE variation is 3.3\% on ETTm1 and 1.3\% on ETTh1 across the
full range tested. This is descriptive single-run sensitivity evidence rather than a
formal hypothesis test. This suggests using $P=16$ as a reasonable default while
acknowledging that
dataset-specific tuning (larger patches for higher-frequency data, smaller for
lower-frequency) may yield marginal improvements.

\subsection{Summary of Empirical Findings}
The evidence supports a bounded conclusion. Weather is the most favorable supporting
dataset. ETTm2 is also somewhat favorable, though the capacity-matched comparison is mixed at
$H=336$ for \modelvtwo. On ETTh2, the same anti-capacity pattern holds within the
PatchTST family, but TiDE has lower error in the broader comparison at longer horizons. On
Illness, \modelvzero\ has the lowest MSE in the standard single-seed benchmark, but the aligned
capacity-control follow-up shows lower values for PatchTST$_\text{cap}$ at $H=24$ and $H=48$ and
PatchTST$_\text{base}$ at $H=60$. ETTm1, ETTh1, and Exchange are neutral to negative.
The K-depth ablation suggests that $K=1$ is a reasonable efficiency-oriented default in the tested setting.

\begin{table*}[t]
\centering
\caption{Claim--evidence--boundary summary for the manuscript's main conclusions.}
\label{tab:claims}
\small
\begin{tabular}{@{}p{0.24\textwidth}p{0.38\textwidth}p{0.30\textwidth}@{}}
\toprule
Claim & Evidence & Boundary \\
\midrule
Fusion may help in selected regimes & Main benchmark plus Table~\ref{tab:capmatch}
show the most favorable evidence on Weather; ETTh2 and ETTm2 provide additional
PatchTST-family evidence; width-matched, depth-matched ($H=192$), and head-capacity
(LargeHead) controls all support the same direction on Weather and ETTh2. & Not universal; Exchange and ETTm1 remain weak or
neutral cases. \\
Some gains are not fully explained by parameter count & Width-matched controls in Table~\ref{tab:capmatch}
across all three horizons; depth-matched PatchTST$_\text{depth}$ ($\sim$3.2M) at $H=192$;
PatchTST$_\text{LargeHead}$ head-capacity control (3.45M $>$ \modelvzero\ 3.26M, single-seed).
& Width and depth are one-dimensional capacity axes; full FLOP-matched or
exhaustive multi-seed capacity designs remain future work. \\
Illness is more consistent with low-data capacity effects & Table~\ref{tab:illness_standard}
gives lower values for \modelvzero\ against standard baselines, while Table~\ref{tab:illness_capacity}
shows lower values for PatchTST variants in the aligned follow-up. & Fusion-specific advantage is not
cleanly isolated in this regime. \\
$K=1$ appears sufficient in the tested settings & Table~\ref{tab:kdepth} shows that $K=1$ gives the
lowest MSE in four of six tested settings. & Single-seed follow-up evidence only; no
formal statistical test. \\
Robustness is regime-dependent & Random/block missingness, Gaussian noise, and channel
dropout tested on ETTm2, Weather, Exchange, and Illness (Table~\ref{tab:robustness_summary},
Figure~\ref{fig:robustness}; supplementary details in Appendix~\ref{app:full_results});
PFB competitive on ETTm2 random missingness; DLinear leads on Exchange and Weather
noise. & Test-time perturbations only; no
retraining or imputation module; broader robustness cannot be guaranteed. \\
Mechanistic diagnostics support regime interpretation & CKA similarity 0.42--0.75
before/after BERT refinement on Weather and ETTm2; dataset-regime diagnostics;
Weather per-variable analysis showing gains on thermodynamic channels. & Diagnostic
evidence only; not causal proof. \\
Pooled significance tests support global trend & Wilcoxon signed-rank test ($N=45$,
Holm--Bonferroni + BH-FDR) on five-seed $H=192$ pooled comparisons. & Pooled
test cannot isolate individual dataset effects; per-dataset tests remain
underpowered. \\
\bottomrule
\end{tabular}
\end{table*}

\subsection{Regime Diagnostics and Representation Evidence}
\label{sec:diagnostics}
To provide mechanistic context for why post-encoder refinement helps in some regimes
but not others, we report three complementary diagnostics: (i) Centered Kernel
Alignment (CKA) similarity before and after BERT refinement, (ii) dataset-level
temporal structure characterization, and (iii) Weather channel-level gain analysis.
These diagnostics are observational and should not be interpreted as causal proofs.

\textbf{CKA similarity before and after refinement.}
Table~\ref{tab:cka} reports CKA similarity between patch representations before and
after the BERT refinement stage on Weather and ETTm2 ($H=192$, seed=2021). A CKA
of 1 would indicate a pass-through (no refinement); lower values indicate more
substantial transformation.

\begin{table}[t]
\centering
\caption{CKA similarity between pre- and post-BERT-refinement representations
(Weather and ETTm2, $H=192$, seed=2021). Values below 1 indicate that the
refinement stage is not a simple pass-through.}
\label{tab:cka}
\small
\begin{tabular}{@{}llc@{}}
\toprule
Dataset & Model & CKA (before/after) \\
\midrule
Weather & \modelvzero & 0.751 \\
Weather & \modelvtwo  & 0.659 \\
ETTm2   & \modelvzero & 0.633 \\
ETTm2   & \modelvtwo  & 0.424 \\
\bottomrule
\end{tabular}
\end{table}

CKA values range from 0.42 to 0.75, indicating that the BERT refinement stage
produces meaningfully transformed representations rather than functioning as a
near-identity mapping. The greater transformation in ETTm2 (especially for
\modelvtwo, CKA=0.424) is consistent with the higher structural regularity of that
dataset (see below). Figure~\ref{fig:cka_similarity} extends this view with
cross-stream comparisons. On ETTm2, PatchTST stays much closer to the PFB patch
stream than to the refined stream, and the two fusion variants are also more aligned
at the patch level than after refinement, supporting the interpretation that the
refinement stage is a major source of representational divergence.

\begin{figure*}[!t]
\centering
\includegraphics[width=\textwidth]{figures-300426/Fig7_cka_similarity.png}
\caption{Expanded CKA similarity view on Weather and ETTm2 ($H=192$, seed=2021),
including before/after refinement within each PFB variant and cross-model stream
comparisons. Higher values indicate stronger representational alignment. The patch
streams remain relatively aligned with PatchTST, while the refined streams diverge
more clearly.}
\label{fig:cka_similarity}
\end{figure*}

\textbf{Dataset-regime diagnostics.}
Table~\ref{tab:regime} characterizes each dataset by lag-1 autocorrelation,
normalized spectral entropy, and dominant spectral period, and assigns a qualitative
regime label based on these statistics.

\begin{table*}[t]
\centering
\caption{Dataset temporal-structure diagnostics. Lag-1 AC and spectral entropy are
computed from the training split of each dataset. Dominant period is the
highest-power spectral peak. Regime labels (Support/Mixed/Failure) reflect
whether PFB gains are generally positive, mixed, or absent based on the
main benchmark results.}
\label{tab:regime}
\small
\begin{tabular}{@{}lcccl@{}}
\toprule
Dataset & Lag-1 AC & Spectral entropy & Dominant period & Regime \\
\midrule
ETTh1    & 0.937 & 0.325 & 24  & Mixed   \\
ETTh2    & 0.969 & 0.157 & 24  & Support \\
ETTm1    & 0.978 & 0.224 & 96  & Mixed   \\
ETTm2    & 0.988 & 0.099 & 96  & Support \\
Exchange & 0.999 & 0.028 & $\sim$1770 & Failure \\
Weather  & 0.892 & 0.246 & 144 & Support \\
Illness  & 0.968 & 0.398 & 52  & Support (capacity-sensitive) \\
\bottomrule
\end{tabular}
\end{table*}

The failure case (Exchange) exhibits near-unit-root dynamics (lag-1 AC=0.999) and
near-zero spectral entropy, consistent with a near-random-walk regime in which
patch-level structural patterns are minimal. The support cases (ETTh2, ETTm2,
Weather, Illness) have richer periodicity and more varied spectral content, consistent
with the intuition that post-encoder refinement is more useful when patch-level
temporal patterns are rich enough to be further refined.

\textbf{Weather channel-level analysis.}
For Weather ($H=192$, seed=2021), channel-level MSE comparison between \modelvzero\
and PatchTST shows gains concentrated on thermodynamic channels: T, Tpot, VPmax,
Tlog, rh, and rho show \modelvzero\ advantages, while the rain and raining channels
show no consistent gain. This per-variable pattern is consistent with the hypothesis
that post-encoder refinement captures structured multivariate dynamics better in
physically correlated channels.


% =========================
\section{Discussion}\label{sec:discussion}
% =========================

The results position \model\ as a targeted extension of patch-based forecasting
rather than as a uniformly superior architecture.

\subsection{When Does Additional Refinement Appear to Help?}
The general pattern is that refinement-and-fusion is most likely to help when the baseline
patch representation is strong but not fully sufficient, and when the dataset has
richer temporal structure. Weather is the most favorable example. ETTm2 is also
somewhat favorable, although the capacity-matched comparison is mixed at $H=336$ because
PatchTST$_\text{cap}$ slightly exceeds \modelvtwo. ETTh2 is consistent with the same interpretation
within the PatchTST family, though TiDE has the lowest error there overall. In
these settings, capacity scaling alone does not generally reproduce the fusion gains,
which is consistent with an architectural effect rather than a pure width
effect.

In contrast, Exchange and ETTm1 are settings where simpler or shallower inductive
biases already capture most of the signal structure, and additional refinement adds
limited practical value.

\subsection{The Illness Finding and Its Updated Interpretation}
The original paper positioned Illness as a central result. The capacity-matched
control experiment revises this interpretation. \modelvzero\ still outperforms the
standard single-seed baselines on Illness, but the aligned multi-seed
capacity-control table changes the picture: PatchTST$_\text{cap}$ has the lowest mean MSE at
$H=24$ and $H=48$, while PatchTST$_\text{base}$ has the lowest mean MSE at $H=60$.

A more cautious framing is: Illness shows that patch-based models with
additional capacity benefit the low-data regime, but the specific contribution of
fusion versus raw capacity is not cleanly separable. Exchange and Illness together
define the boundaries of the architectural claim: Exchange is a case where fusion does not help,
and Illness is a case where the fusion-specific effect is most uncertain.

\subsection{Interpreting the Two Fusion Variants}
Neither variant dominates. When direct concatenation performs better, explicit
retention of both streams may be helpful. When projection performs better, learned
compression and reweighting of fused features may also be helpful. This suggests treating
fusion design as a meaningful architectural component. The K-depth ablation adds a
further dimension: the choice of $K$ is less important than the choice of fusion
mechanism, as performance is largely flat over $K \in \{1,2,3\}$.

\subsection{Relationship Between Accuracy and Complexity}
The capacity-matched control experiment is intended to address the central confound
identified in the original paper. On Weather and ETTh2, both fusion variants
outperform PatchTST at matched parameter counts across all tested horizons. On ETTm2,
\modelvzero\ remains the lowest-MSE fusion variant, but \modelvtwo\ is not uniformly better than
PatchTST$_\text{cap}$ at $H=336$. Taken together, these results are consistent with an
architectural effect in the more favorable regimes while also showing that the
anti-capacity argument is not universal across every variant and horizon.

The remaining complexity-related caveat is that these controls cover selected
capacity axes --- width matching, one depth-matched $H=192$ setting, and one
single-seed LargeHead diagnostic --- rather than a full FLOP-matched or
exhaustive architecture-matching study.

\subsection{Controlled Efficiency Benchmark}
Table~\ref{tab:efficiency_compact} reports a controlled efficiency benchmark on
Weather $H=192$ using three repeated runs on an NVIDIA RTX 2080~Ti, with identical
hardware and data conditions for all models. Training time, inference latency, GPU
memory, and MACs are the median of three runs. FLOPs are computed analytically.
Figure~\ref{fig:efficiency_pareto}
provides a complementary visual summary of the accuracy--cost trade-off.

DLinear remains by far the fastest and most memory-efficient model (321\,s training,
0.33\,ms inference latency). Among patch-based models, PFBv0 trains in
approximately the same wall-clock time as PatchTST$_\text{cap}$ (1{,}375\,s
vs.\ 1{,}378\,s) but uses 30\% fewer MACs (1.09G vs.\ 1.56G) and requires
substantially less inference peak memory (72\,MB vs.\ 193\,MB). PFBv0 is not faster
than the lightweight PatchTST (1{,}018\,s) or DLinear; the efficiency gains are
relative to the capacity-matched control, not absolute.

\begin{table*}[t]
\centering
\caption{Controlled efficiency benchmark on Weather $H=192$ (median of 3 repeated
runs, NVIDIA RTX~2080~Ti). All models trained and evaluated under identical
hardware and data conditions. MACs computed with \texttt{torchinfo} for a single
forward pass over one batch; FLOPs $\approx 2\times$MACs.
\textdagger\ PatchTST$_\text{cap}$ uses the capacity-matched configuration
(3.28M params at H=192).}
\label{tab:efficiency_compact}
\small
\resizebox{\textwidth}{!}{%
\begin{tabular}{@{}lcccccc@{}}
\toprule
Model & Params (H=192) & Train median (s) & Peak GPU train (MB) & Infer.\ median (ms) & Infer.\ peak GPU (MB) & MACs \\
\midrule
\dlinear    & 0.13M & 321   & 1{,}966 & 0.33 & 10  & 2.72M \\
\patchtst   & 1.63M & 1{,}018 & 2{,}812 & 3.25 & 62  & 546.87M \\
PatchTST$_\text{cap}$\textdagger & 3.28M & 1{,}378 & 3{,}107 & 6.64 & 193 & 1.56G \\
\modelvzero & 3.26M & 1{,}375 & 3{,}205 & 5.87 & 72  & 1.09G \\
\modelvtwo  & 2.46M & 1{,}461 & 1{,}994 & 6.75 & 69  & 1.27G \\
\bottomrule
\end{tabular}
}
\end{table*}

\begin{figure}[t]
\centering
\includegraphics[width=0.98\linewidth]{figures-300426/Fig4_efficiency_pareto.png}
\caption{Accuracy--efficiency trade-off on Weather $H=192$ using the controlled
benchmark values from Table~\ref{tab:efficiency_compact}. Aggregate MSE is a
descriptive cross-setting summary from the main benchmark archive. MACs are
from the controlled profiling run.}
\label{fig:efficiency_pareto}
\end{figure}

\subsection{Limitations}
First, capacity attribution is supported by width-matched PatchTST controls across
horizons, a depth-matched PatchTST control at $H=192$, and a single-seed LargeHead
head-capacity diagnostic. However, full FLOP-matched, multi-seed, and broader
architecture-matching capacity designs remain future work.

Second, multi-seed evaluation is representative rather than exhaustive:
multi-seed stability analyses use five seeds (2021--2025); the capacity-matched control
and Illness capacity-control experiments use three runs. Although the study includes
pooled Wilcoxon signed-rank tests with Holm--Bonferroni and BH-FDR correction,
per-configuration statistical testing remains underpowered with five seeds.
Dataset-level conclusions therefore rely on mean performance, standard deviation,
effect direction, and capacity-control consistency rather than per-dataset
significance claims.

Third, robustness probes cover random and block missingness, Gaussian additive noise,
and channel dropout across multiple datasets, but all are test-time perturbations
applied without retraining or imputation. They should be treated as diagnostic stress
tests rather than a guarantee of broad deployment robustness.

Fourth, the study focuses on supervised point forecasting. Extension to probabilistic
forecasting, transfer settings, and foundation-model pretraining are left for future
work.

Fifth, no per-baseline hyperparameter tuning was performed. All models were trained
under a single shared protocol (patch models LR=$10^{-4}$; DLinear LR=$10^{-2}$;
iTransformer, TiDE, and TimeXer LR=$10^{-3}$), which is a fair but not individually
tuned setting; baselines whose canonical implementations use different optimizers,
learning rates, or schedules may be disadvantaged relative to their reported figures
in prior papers.

Sixth, the pooled Wilcoxon tests aggregate Illness horizons with standard horizons,
which may inflate power due to scale differences. Per-configuration tests remain
underpowered ($N=5$), so dataset-level conclusions rely on directional evidence,
multi-seed variability, and capacity controls rather than per-dataset significance.

Seventh, this study does not cover ECL (Electricity) or Traffic benchmarks, and
$H=720$ is outside the tested horizon range. The 21-setting protocol was chosen for
cross-dataset split consistency; extending to these additional settings remains
future work.

% =========================
\section{Conclusion}
% =========================

This paper introduced \model, a patch-based long-horizon forecasting architecture
that augments a PatchTST-style backbone with a compact post-encoder refinement stage
and explicit fusion between backbone and refined representations. The central objective
was to test a focused architectural question: whether additional refinement of patch
representations can improve forecasting quality without abandoning the efficiency
advantages of patch-based tokenization.

The empirical evidence suggests a structured, regime-dependent
answer. Weather provides the most favorable evidence overall: under the follow-up
capacity-control protocol, both fusion variants outperform matched PatchTST at all
three horizons, while the main benchmark shows relatively favorable results for
\modelvtwo\ and competitive performance for \modelvzero. ETTm2
remains somewhat favorable but mixed at $H=336$ in the capacity-matched comparison, and ETTh2
is broadly consistent with the same interpretation within the PatchTST family even though TiDE is the
lowest-error baseline at longer horizons. Gains are not universal: Exchange and
ETTm1 indicate that simpler baselines remain competitive in settings where the
temporal structure does not benefit from additional refinement. On Illness,
\modelvzero\ outperforms standard baselines in the main single-seed benchmark, but
the aligned capacity-control multi-seed follow-up favors PatchTST$_\text{cap}$ at
$H=24$ and $H=48$ and PatchTST$_\text{base}$ at $H=60$, indicating that added
patch-based capacity rather than fusion is the more cautious interpretation in that
low-data regime.

The K-depth ablation shows that a single refinement block captures most of the
architectural benefit in the tested single-seed setting. The Pointwise+Refine
ablation shows that deeper refinement alone is not consistently competitive without
patch tokenization and explicit fusion. Neither fusion variant dominates the other consistently,
suggesting that fusion design is a meaningful design choice.

Future work should prioritize more systematic parameter-matched attribution studies,
broader fusion and refinement-depth ablations, wider multi-seed coverage, richer
robustness protocols, and extension to probabilistic forecasting and transfer settings.

\FloatBarrier

\begin{thebibliography}{99}

\bibitem{zhou2021informer}
H.~Zhou, S.~Zhang, J.~Peng, S.~Zhang, J.~Li, H.~Xiong, and W.~Zhang.
\newblock Informer: Beyond Efficient Transformer for Long Sequence Time-Series
Forecasting.
\newblock In \emph{Proc.\ AAAI}, 2021.

\bibitem{wu2021autoformer}
H.~Wu, J.~Xu, J.~Wang, and M.~Long.
\newblock Autoformer: Decomposition Transformers with Auto-Correlation for Long-Term
Series Forecasting.
\newblock In \emph{Advances in Neural Information Processing Systems (NeurIPS)}, 2021.

\bibitem{kim2022revin}
T.~Kim, J.~Kim, Y.~Tae, C.~Park, J.-H. Choi, and J.~Choo.
\newblock Reversible Instance Normalization for Accurate Time-Series Forecasting
against Distribution Shift.
\newblock In \emph{Proc.\ ICLR}, 2022.

\bibitem{vaswani2017attention}
A.~Vaswani, N.~Shazeer, N.~Parmar, J.~Uszkoreit, L.~Jones, A.~N. Gomez,
\L{}.~Kaiser, and I.~Polosukhin.
\newblock Attention Is All You Need.
\newblock In \emph{Advances in Neural Information Processing Systems (NeurIPS)}, 2017.

\bibitem{zeng2023transformers}
A.~Zeng, M.~Chen, L.~Zhang, and Q.~Xu.
\newblock Are Transformers Effective for Time Series Forecasting?
\newblock In \emph{Proc.\ AAAI}, 2023.

\bibitem{nie2023patchtst}
Y.~Nie, N.~H. Nguyen, P.~Sinthong, and J.~Kalagnanam.
\newblock A Time Series is Worth 64 Words: Long-term Forecasting with Transformers.
\newblock In \emph{Proc.\ ICLR}, 2023.

\bibitem{liu2024itransformer}
Y.~Liu, T.~Hu, H.~Wu, H.~Zhang, S.~Hong, L.~Zheng, C.~Guo, \emph{et~al.}
\newblock iTransformer: Inverted Transformers are Effective for Time Series
Forecasting.
\newblock In \emph{Proc.\ ICLR}, 2024.

\bibitem{wang2024timemixer}
S.~Wang, H.~Wu, X.~Shi, T.~Hu, H.~Luo, L.~Ma, J.~Y. Zhang, and J.~Zhou.
\newblock TimeMixer: Decomposable Multiscale Mixing for Time Series Forecasting.
\newblock In \emph{Proc.\ ICLR}, 2024.

\bibitem{das2023tide}
A.~Das, W.~Kong, A.~Leach, S.~Mathur, R.~Sen, and R.~Yu.
\newblock Long-term Forecasting with TiDE: Time-series Dense Encoder.
\newblock \emph{Transactions on Machine Learning Research (TMLR)}, 2023.

\bibitem{wang2024timexer}
X.~Wang, Y.~Lv, G.~Yin, J.~Liu, \emph{et~al.}
\newblock TimeXer: Empowering Transformers for Time Series Forecasting with Exogenous
Variables.
\newblock \emph{arXiv preprint arXiv:2402.19072}, 2024.

\bibitem{selfattention2024tsf}
Y.~Dong, M.~Yan, Y.~Jiang, \emph{et~al.}
\newblock Are Self-Attentions Effective for Time Series Forecasting?
\newblock In \emph{Advances in Neural Information Processing Systems (NeurIPS)}, 2024.

\bibitem{closerlook2024tsf}
Z.~Li, X.~Liu, \emph{et~al.}
\newblock A Closer Look at Transformers for Time Series Forecasting.
\newblock \emph{OpenReview preprint}, 2024.

\bibitem{ansari2024chronos}
A.~Ansari, O.~Stella, C.~Turkmen, X.~Zhang, P.~Mergenthaler, \emph{et~al.}
\newblock Chronos: Learning the Language of Time Series.
\newblock \emph{arXiv preprint arXiv:2403.07815}, 2024.

\bibitem{timemixerpp2025}
S.~Wang, H.~Wu, \emph{et~al.}
\newblock TimeMixer++: A General Time Series Pattern Machine for Universal
Time-Series Analysis.
\newblock \emph{OpenReview preprint}, 2025.

\bibitem{devlin2019bert}
J.~Devlin, M.-W. Chang, K.~Lee, and K.~Toutanova.
\newblock BERT: Pre-training of Deep Bidirectional Transformers for Language
Understanding.
\newblock In \emph{Proc.\ NAACL-HLT}, 2019.

\bibitem{cheng2023timemae}
M.~Cheng, Q.~Liu, Z.~Liu, H.~Zhang, R.~Zhang, and E.~Chen.
\newblock TimeMAE: Self-Supervised Representations of Time Series with Decoupled
Masked Autoencoders.
\newblock \emph{arXiv preprint arXiv:2303.00320}, 2023.

\bibitem{lai2018lstnet}
G.~Lai, W.-C. Chang, Y.~Yang, and H.~Liu.
\newblock Modeling Long- and Short-Term Temporal Patterns with Deep Neural Networks.
\newblock In \emph{Proc.\ ACM SIGIR}, 2018.

\bibitem{ahmad2024icamcs_patchbert}
H.~Ahmad, S.~K.~Mortazavi, S.~Elbahnasawi, F.~Al~Machot, and K.~Kyamakya.
\newblock Enhanced Time Series Forecasting: Integrating PatchTST with BERT Layers.
\newblock In \emph{Proc.\ 2024 International Conference on Applied Mathematics and
Computer Science (ICAMCS)}. IEEE, 2024.

\end{thebibliography}

\appendices
\section{Supplementary Experimental Results}\label{app:full_results}

For readability, this appendix is organized into three blocks: (i) stability analyses and
(ii) robustness/sensitivity analyses, and (iii) supplementary diagnostic/protocol-
extension views, followed by artifact availability information.

\subsection{Stability Analyses}

\subsubsection{Extended Multi-Seed Summary}
Table~\ref{tab:app_multiseed_mse} reports MSE mean$\pm$std over five seeds (2021--2025) for
representative $H=192$ settings across six datasets.

\begin{table*}[t]
\centering
\caption{Representative multi-seed MSE summary (mean$\pm$std over five seeds: 2021--2025, lower is better)
for $H=192$.}
\label{tab:app_multiseed_mse}
{\small
\setlength{\tabcolsep}{12pt}
\renewcommand{\arraystretch}{1.15}
\begin{tabular}{@{}lccc@{}}
\toprule
Dataset & DLinear & PatchTST & \modelvzero \\
\midrule
ETTm1    & \textbf{0.3351 $\pm$ 0.0020} & 0.3364 $\pm$ 0.0041 & 0.3395 $\pm$ 0.0023 \\
ETTm2    & 0.2308 $\pm$ 0.0051          & 0.2370 $\pm$ 0.0055 & \textbf{0.2287 $\pm$ 0.0042} \\
ETTh1    & \textbf{0.4074 $\pm$ 0.0023} & 0.4321 $\pm$ 0.0014 & 0.4196 $\pm$ 0.0101 \\
ETTh2    & 0.3945 $\pm$ 0.0072          & 0.3824 $\pm$ 0.0071 & \textbf{0.3740 $\pm$ 0.0083} \\
Exchange & \textbf{0.1779 $\pm$ 0.0056} & 0.1921 $\pm$ 0.0052 & 0.2244 $\pm$ 0.0084 \\
Weather  & 0.2168 $\pm$ 0.0020          & 0.1976 $\pm$ 0.0016 & \textbf{0.1964 $\pm$ 0.0019} \\
\bottomrule
\end{tabular}
}
\end{table*}

Figure~\ref{fig:app_full_errorbars} complements the table with an error-bar view
of the same representative settings.

\begin{figure*}[!t]
\centering
\includegraphics[width=\textwidth]{figures-300426/A1_full_errorbars.png}
\caption{Companion multi-seed error-bar view for the representative $H=192$
settings in Table~\ref{tab:app_multiseed_mse}. Means are shown with variability
bars over the five-seed runs.}
\label{fig:app_full_errorbars}
\end{figure*}

Figure~\ref{fig:app_forecast_dispersion} provides a qualitative trajectory-level
view of seed-to-seed forecast dispersion. It is included only as a complement to the
numerical multi-seed tables and should not be interpreted as calibrated uncertainty
quantification.

\begin{figure*}[!t]
\centering
\includegraphics[width=\textwidth]{figures-300426/Fig_forecast_uncertainty.png}
\caption{Illustrative trajectory-level forecast-dispersion examples from the
repeated-run setting. The visualization summarizes empirical seed dispersion around
the point forecasts and is included only as a qualitative complement to the
multi-seed tables, not as calibrated uncertainty quantification.}
\label{fig:app_forecast_dispersion}
\end{figure*}

Figure~\ref{fig:app_significance} provides a compact visual summary of the pooled
and horizon-specific Wilcoxon outcomes discussed in Section~\ref{sec:wilcoxon}.

\begin{figure*}[!t]
\centering
\includegraphics[width=\textwidth]{figures-300426/A2_significance.png}
\caption{Supplementary visual summary of the pooled Wilcoxon signed-rank outcomes
reported in Section~\ref{sec:wilcoxon}. The display condenses the corrected and
uncorrected comparisons into a single view, distinguishing Holm-significant,
BH-FDR-only, and directional/non-significant results.}
\label{fig:app_significance}
\end{figure*}

\subsubsection{Illness Multi-Seed Stability}
Table~\ref{tab:app_illness_multiseed} reports Illness multi-seed results from the
aligned 3-seed capacity-control follow-up (mean$\pm$std, LR=$10^{-4}$). The main observation is
that PatchTST$_\text{cap}$ has lower mean MSE and lower variability at $H=24$ and
$H=48$ (std 0.021 and 0.003) compared to \modelvzero\ (std 0.169 and 0.024). At
$H=60$, PatchTST$_\text{base}$ achieves the lowest mean MSE and the lowest variance
among the top-performing models.

\begin{table*}[t]
\centering
\caption{Illness multi-seed stability (mean $\pm$ std over 3 seeds,
follow-up protocol LR=$10^{-4}$; lower is better).}
\label{tab:app_illness_multiseed}
{\small
\setlength{\tabcolsep}{10pt}
\renewcommand{\arraystretch}{1.12}
\begin{tabular}{@{}llcc@{}}
\toprule
Horizon & Model & MSE (mean$\pm$std) & MAE (mean$\pm$std) \\
\midrule
\multirow{4}{*}{$H=24$}
& PatchTST$_\text{base}$ & 2.076 $\pm$ 0.132 & 0.947 $\pm$ 0.071 \\
& PatchTST$_\text{cap}$  & \textbf{1.899 $\pm$ 0.021} & \textbf{0.862 $\pm$ 0.008} \\
& \modelvzero             & 2.061 $\pm$ 0.169 & 0.919 $\pm$ 0.065 \\
& \modelvtwo              & 2.200 $\pm$ 0.148 & 0.982 $\pm$ 0.060 \\
\midrule
\multirow{4}{*}{$H=48$}
& PatchTST$_\text{base}$ & 2.078 $\pm$ 0.154 & 0.966 $\pm$ 0.064 \\
& PatchTST$_\text{cap}$  & \textbf{1.892 $\pm$ 0.003} & \textbf{0.886 $\pm$ 0.004} \\
& \modelvzero             & \underline{1.911 $\pm$ 0.024} & \underline{0.898 $\pm$ 0.007} \\
& \modelvtwo              & 2.058 $\pm$ 0.061 & 0.940 $\pm$ 0.018 \\
\midrule
\multirow{4}{*}{$H=60$}
& PatchTST$_\text{base}$ & \textbf{2.000 $\pm$ 0.019} & \textbf{0.928 $\pm$ 0.004} \\
& PatchTST$_\text{cap}$  & 2.097 $\pm$ 0.099 & 0.965 $\pm$ 0.043 \\
& \modelvzero             & 2.036 $\pm$ 0.071 & 0.959 $\pm$ 0.042 \\
& \modelvtwo              & 2.072 $\pm$ 0.016 & 0.952 $\pm$ 0.005 \\
\bottomrule
\end{tabular}
}
\end{table*}

\FloatBarrier

\subsection{Robustness and Sensitivity Analyses}

\subsubsection{Original ETTm1 Missing-Input Robustness Case Study}
For reference, Table~\ref{tab:app_robustness_full} reports the full numeric values
for the original ETTm1 $H=192$ random-missingness case study discussed in
Section~\ref{sec:results}. This is a single-dataset single-corruption case study. The
broader robustness summary across multiple datasets and corruption types is in
Table~\ref{tab:robustness_summary}. Figure~\ref{fig:app_robustness_matrix} provides a
compact matrix visualization of the ETTm1 case study.
\begin{table}[!htbp]
\centering
\caption{Missing-input robustness (MSE, ETTm1 $H=192$). Lower is better.}
\label{tab:app_robustness_full}
{\small
\setlength{\tabcolsep}{4pt}
\renewcommand{\arraystretch}{1.12}
\begin{tabular}{@{}lccccc@{}}
\toprule
Model & $\rho=0.0$ & $\rho=0.1$ & $\rho=0.2$ & $\rho=0.3$ & \shortstack{Deg.\\(30\%)} \\
\midrule
PFBv0    & 0.3426 & 0.3584 & 0.3839 & 0.4191 & +22.3\% \\
PatchTST & 0.3429 & 0.3594 & 0.3833 & 0.4184 & +22.0\% \\
DLinear  & 0.3345 & 0.3515 & 0.3826 & 0.4261 & +27.4\% \\
\bottomrule
\end{tabular}
}
\end{table}

\begin{figure*}[!t]
\centering
\includegraphics[width=\textwidth]{figures-300426/A7_robustness_matrix.png}
\caption{Supplementary robustness matrix view for the missing-input case study,
providing a compact comparison across models and masking levels for the ETTm1
$H=192$ setting. Lower is better.}
\label{fig:app_robustness_matrix}
\end{figure*}

Figure~\ref{fig:app_gaussian_channel_robustness} supplements the main robustness
summary with Gaussian-noise and channel-dropout views for ETTm2 and Weather.

\begin{figure*}[!t]
\centering
\includegraphics[width=\textwidth]{figures-300426/A8_gaussian_channel_robustness.png}
\caption{Supplementary Gaussian-noise and channel-dropout robustness views for
ETTm2 and Weather at $H \in \{96,192\}$ under inference-only perturbations. Lower
degradation is better. The plots highlight DLinear's strong noise robustness on
Weather and the severe effect of channel dropout on ETTm2 for all models.}
\label{fig:app_gaussian_channel_robustness}
\end{figure*}

\subsubsection{Patch-Length Sensitivity}
Table~\ref{tab:app_patch_sensitivity} restates the patch-length sensitivity summary discussed in the main text.
\begin{table}[!htbp]
\centering
\caption{Patch-length sensitivity for \modelvzero\ (lower MSE is better).
PatchTST results are excluded as invariance was observed (likely checkpoint reuse).}
\label{tab:app_patch_sensitivity}
{\small
\setlength{\tabcolsep}{6pt}
\renewcommand{\arraystretch}{1.12}
\begin{tabular}{@{}lcccc@{}}
\toprule
Dataset & Best $P$ & Best MSE & MSE Range & Relative Range \\
\midrule
ETTm1 & 32 & 0.3293 & 0.0109 & 3.26\% \\
ETTh1 & 8  & 0.4069 & 0.0053 & 1.29\% \\
\bottomrule
\end{tabular}
}
\end{table}

\subsection{Supplementary Diagnostic and Protocol-Extension Views}

\subsubsection{Error-by-Step and Channel-Level Views}
Figure~\ref{fig:app_error_by_step} expands the Weather/ETTm2 discussion by showing
how error evolves across the forecast horizon. Figure~\ref{fig:app_per_variable_weather}
complements the Weather discussion with a channel-wise breakdown at $H=192$.

\begin{figure*}[!t]
\centering
\includegraphics[width=\textwidth]{figures-300426/A4_error_by_step.png}
\caption{Per-step MSE across the forecast horizon for Weather and ETTm2 at
$H=192$, using the repeated-run analysis described in the text. Lower is better.
On Weather, the gap between PFBv0 and DLinear becomes more pronounced at longer
lead times; on ETTm2, DLinear is strongest at the earliest steps while PFBv0
becomes more competitive later in the horizon.}
\label{fig:app_error_by_step}
\end{figure*}

\begin{figure*}[!t]
\centering
\includegraphics[width=\textwidth]{figures-300426/A3_per_variable_weather.png}
\caption{Per-variable Weather MSE comparison at $H=192$ (seed 2021) for DLinear,
PatchTST, and \modelvzero. Gains are concentrated on thermodynamic channels such as
T, Tpot, VPmax, Tlog, rh, and rho, whereas rain-related channels show little or no
advantage.}
\label{fig:app_per_variable_weather}
\end{figure*}

\subsubsection{Unified-LR Bridging View}
As noted in Section~\ref{sec:experimental_setup}, the figure below extends the
single-seed unified-LR protocol to include supplementary $H=720$ runs. It is shown
only as descriptive bridging evidence and is not part of the locked 21-setting
benchmark.

\begin{figure*}[!t]
\centering
\includegraphics[width=\textwidth]{figures-300426/A6_bridging_lr.png}
\caption{Single-seed unified-LR bridging view across datasets and horizons,
including the supplementary $H=720$ runs. This figure is descriptive bridging
evidence only and is not part of the locked 21-setting benchmark.}
\label{fig:app_bridging_lr}
\end{figure*}

% \subsection{Artifact Availability}
% Complete machine-readable outputs are provided under
% \texttt{patchfusionbert-results/24-03-2026} in the \texttt{csv} and \texttt{md}
% subdirectories, including merged result tables, aggregated statistics, and narrative
% reports. These files support independent re-analysis beyond the manuscript tables.

\clearpage

\section*{Data Availability}
All datasets used in this study (ETTh1, ETTh2, ETTm1, ETTm2, Exchange, Weather,
Illness) are publicly available through their respective benchmark repositories
under their original licenses.
Code for reproducing the experiments is available at
\url{https://github.com/Hussein-Ahmad-Ahmad/patchfusionbert-forecasting},
tagged as \texttt{v1.0-submission} (commit \texttt{e2b5fbb}).
The repository includes: (i)~training scripts and per-dataset configuration files
for the main 21-setting grid; (ii)~scripts for the capacity-matched control,
K-depth ablation, and multi-seed stability runs; (iii)~\texttt{requirements.txt}
and \texttt{environment.yml} specifying all package dependencies; (iv)~the locked
result archive (\texttt{results/locked\_snapshot/}) from which all manuscript
tables are assembled; (v)~a \texttt{README.md} with a directory map and
reproduction instructions.
\textbf{Hardware and software:} NVIDIA RTX 2080 Ti (11\,GB), CUDA 11.8,
PyTorch 2.0.1, Python 3.10, Ubuntu 20.04.
Approximate per-configuration runtime: 1--4 GPU-hours (main grid).
Full-grid reproduction is approximately 80 GPU-hours.


% --- BIOGRAPHIES ---
\begin{IEEEbiography}[{\safeincludegraphics[width=1in,height=1.25in,clip,keepaspectratio]{figures/authors/hussein.jpg}}]{Hussein Ahmad}
received the B.Sc.
degree in Mathematics from the Lebanese University,
Lebanon, and the M.Sc. degree in Applied
Mathematics from the University of Debrecen,
Hungary. He is currently pursuing the Ph.D. degree
in Information and Communications Engineering
within the Doctoral Programme in Technical
Sciences at the University of Klagenfurt,
Austria.
His research interests include trustworthy and
efficient spatio-temporal forecasting, machine learning, explainable AI, predictive
modelling, data analytics, and time-series forecasting.
\end{IEEEbiography}

\begin{IEEEbiography}[{\safeincludegraphics[width=1in,height=1.25in,clip,keepaspectratio]{figures/authors/kasra.jpg}}]{Seyyed Kasra Mortazavi}
is a Ph.D. candidate
in the Department of Smart Systems Technologies
at the University of Klagenfurt, Austria.
His research interests include time-series forecasting,
intelligent energy management, machinelearning-
based resource allocation, 5G network
slicing, and autonomous vehicular systems.
He received the M.Sc. degree in Electronics
Engineering from Amirkabir University of Technology
(Tehran Polytechnic) and the B.Sc. degree
in Electronics Technology from Shiraz University. He has over seventeen
years of industrial and research experience in telecommunications, IoT, and
smart-energy systems, and is a member of the Young Inventors Organization
of Iran.
\end{IEEEbiography}

\begin{IEEEbiography}[{\safeincludegraphics[width=1in,height=1.25in,clip,keepaspectratio]{figures/authors/taha.jpg}}]{Taha Benarbia}
received the Ph.D. degree in electrical engineering from the University of Cergy--Pontoise (UCP), France in 2013. He is currently an Associate Professor at faculty of maintenance and industrial security at University of Oran 2 Mohamed Ben Ahmed, Oran. In 2017, he joined the Institute for Transport Planning and Systems IVT, ETH Zürich, Switzerland, as a Postdoctoral. His main research interests include shared mobility systems, connected and automated vehicles, drone delivery systems modelling, stochastic and discrete event systems optimization, performance evaluation of dynamic systems, machine learning and neuro computing.
\end{IEEEbiography}

\begin{IEEEbiography}[{\safeincludegraphics[width=1in,height=1.25in,clip,keepaspectratio]{figures/authors/fadi.jpg}}]{Fadi Al Machot}
(Member, IEEE) received the German Diploma degree in computer science from the University of Potsdam, in 2010, the Ph.D. degree in computer science from Klagenfurt University, Austria, in 2013, and the Habilitation degree in applied computer science from the University of Lübeck, Germany, in 2020. He is currently an Associate Professor in machine learning (ML) with Norwegian University of Life Sciences (NMBU), Norway. His work has been patented and published in peer-reviewed international conferences and journals. His research interests include deep learning, neural-symbolic learning, video understanding, cognitive modeling, and zero/few-shot learning. He is an active reviewer of well-known journals.
\end{IEEEbiography}

\begin{IEEEbiography}[{\safeincludegraphics[width=1in,height=1.25in,clip,keepaspectratio]{figures/authors/kya.jpg}}]{Kyandoghere Kyamakya}
is currently a Full Professor of transportation informatics and the Deputy Director of the Institute of Smart Systems Technologies, University of Klagenfurt, Austria. He is actively conducting research on modeling, simulation, and test-bed evaluations for a series of concepts among others in the context of intelligent transportation systems. In the research addressing transportation systems, a series of fundamental and theoretical tools from the fields of applied mathematics, electronics, and computer science is either extensively exploited or a source of inspiration for innovative solutions and concepts, including nonlinear dynamics, systems science, machine learning/deep learning, nonlinear image processing, and neuro computing. He co-edited more than six books. He has published more than 100 journal articles and more than 100 conference papers.
\end{IEEEbiography}

\vspace{12pt}

\noindent\textbf{CRediT Author Statement}

\noindent Hussein Ahmad: Conceptualization, Methodology, Writing, Graphics, Original draft preparation. Seyyed Kasra Mortazavi: Writing. Taha Benarbia: Writing, Reviewing, Supervising. Fadi Al Machot: Reviewing. Kyandoghere Kyamakya: Reviewing, Supervising.

\EOD

\end{document}
